\documentclass[11pt]{article}

\usepackage{amsmath,amssymb,amsthm,mathtools}
\usepackage[margin=1.1in]{geometry}
\usepackage{microtype}
\usepackage{enumitem}
\usepackage[colorlinks=true,linkcolor=blue,citecolor=blue]{hyperref}
\usepackage{cleveref}

\theoremstyle{plain}
\newtheorem{theorem}{Theorem}[section]
\newtheorem{lemma}[theorem]{Lemma}
\newtheorem{proposition}[theorem]{Proposition}
\newtheorem{corollary}[theorem]{Corollary}
\theoremstyle{definition}
\newtheorem{definition}[theorem]{Definition}
\newtheorem{example}[theorem]{Example}
\newtheorem{remark}[theorem]{Remark}
\newtheorem{problem}[theorem]{Problem}

\DeclareMathOperator{\val}{val}
\DeclareMathOperator{\Opt}{Opt}
\DeclareMathOperator{\opt}{opt}
\DeclareMathOperator{\lfp}{lfp}
\newcommand{\Vmax}{V_{\max}}
\newcommand{\Vmin}{V_{\min}}
\newcommand{\Vavg}{V_{\mathrm{avg}}}
\newcommand{\bbone}{\mathbf{1}}

\title{Simple Stochastic Games:\\
A Self-Contained Development of the Value Problem,\\
and the Exact Location of the Remaining Gap}
\author{}
\date{}

\begin{document}
\maketitle

\begin{abstract}
We give a self-contained, fully proved development of the theory
surrounding the decision problem \textsc{SSG-Value}: given a simple
stochastic game and a start vertex $v_0$, decide whether
$\val(v_0)\ge\tfrac12$. Everything asserted below is proved from first
principles, and every negative statement is witnessed by an explicit
instance verified in exact rational arithmetic. We establish: a
quantitative stopping transformation reducing the problem exactly
(\Cref{thm:stopping-transform}); the complete fixed-point theory of
stopping games (\Cref{sec:stopping}); a structure theorem identifying the
set of optimal Max strategies with a combinatorial subcube
(\Cref{thm:opt-subcube}), together with an explicit witness that the
hypothesis cannot be dropped (\Cref{ex:nonstopping-nonsubcube}); a
polynomial special case --- \emph{arbitrary} SSGs with $O(\log N)$
average vertices (\Cref{thm:few-avg}) --- together with the one-player
boundary (\Cref{thm:one-player}) and a constraint it imposes on any
proposed complexity parameter (\Cref{rem:owner-blind}); a subexponential Las~Vegas algorithm of
expected running time $e^{2\sqrt n}\,\mathrm{poly}(N)$
(\Cref{thm:subexp}), which is the strongest upper bound we obtain; and a
\emph{short improving path} theorem (\Cref{thm:short-path}) showing that
from every strategy some sequence of at most $|\Vmax|$ single strictly
improving switches reaches an optimal strategy. Consequently the entire
difficulty of \textsc{SSG-Value} is concentrated in the \emph{selection}
problem: improving paths of linear length always exist, and no
polynomial-time rule is known to find one --- the most natural rule
provably does not, since all-switches can take more than $|\Vmax|$
iterations and undo its own choices (\Cref{prop:allswitch-overshoot}). We then record four
refutations that close off specific proof strategies --- the
all-switches rule does not dominate partial switches
(\Cref{thm:all-switches-refuted}); the unique-sink orientations induced
by simple stochastic games on the combined two-player cube can be cyclic
(\Cref{thm:cyclic-uso}); value
iteration needs exponentially many steps already on games with no players
at all (\Cref{thm:vi-lower}); and an
improving switch need not reduce the Hamming distance to the optimal set,
so no naive Hamming potential exists (\Cref{thm:hamming-refuted}), with
seven natural polynomial-time selection rules refuted individually
(\Cref{prop:rules-fail}), and a family showing that the guarantee of
\Cref{thm:short-path} is tight --- the switchable vertices worth taking
can be a $1/|\Vmax|$ fraction of all switchable vertices
(\Cref{prop:needle}). We show that \textsc{SSG-Value} is
polynomial-time equivalent to comparing the values of two vertices, and
that $n$ such comparisons --- one per Max vertex, and not even
adaptively chosen --- suffice to determine an optimal strategy (we prove
sufficiency only; no matching lower bound is claimed), indeed that the mere
order of the average-vertex values already determines the whole value
vector (\Cref{thm:order-determines})
(\Cref{thm:compare-equivalence}); the obstruction is therefore that each
of those $n$ bits is itself an instance of the original problem. Finally
we state the remaining gap at full strength (\Cref{sec:gap}).

We emphasise at the outset what this document does \emph{not} contain: it
does not contain a polynomial-time algorithm for \textsc{SSG-Value}, and
it does not claim one. \Cref{sec:gap} states precisely the missing
statement, and proves that it is equivalent in strength to the target.
\end{abstract}

\tableofcontents

\section{Introduction}

\subsection{The problem}

\begin{definition}[Simple stochastic game]\label{def:ssg}
A \emph{simple stochastic game} (SSG) is a finite directed graph $G$ on a
vertex set
\[
  V \;=\; \Vmax \,\cup\, \Vmin \,\cup\, \Vavg \,\cup\, \{t_0,t_1\},
\]
a disjoint union, in which the two \emph{sinks} $t_0,t_1$ have no outgoing
edges and every other vertex has exactly two outgoing edges (parallel
edges and self-loops are permitted). We write $N := |V|$ and
$n := |\Vmax|$, and we denote the two successors of a non-sink vertex $v$
by $v^{(0)}$ and $v^{(1)}$.

A token is placed on a designated start vertex $v_0$ and moved along
edges. At a vertex of $\Vmax$ the edge is chosen by player Max; at a
vertex of $\Vmin$ by player Min; at a vertex of $\Vavg$ by a fair coin,
independently of everything else. Play stops if a sink is reached. Max
wins if and only if the token reaches $t_1$; every other play, including
every infinite play, is a win for Min.
\end{definition}

\begin{definition}[Strategies and values]\label{def:strategies}
A \emph{positional} (equivalently, pure memoryless) strategy for Max is a
map $\sigma\colon \Vmax\to V$ with $\sigma(v)\in\{v^{(0)},v^{(1)}\}$;
positional strategies $\tau$ for Min are defined analogously on $\Vmin$.
We write $\bar\sigma(v)$ for the successor of $v$ \emph{not} chosen by
$\sigma$, and $\sigma[S]$ for the strategy obtained from $\sigma$ by
flipping the choice at every vertex of $S\subseteq\Vmax$; we abbreviate
$\sigma[\{v\}]$ to $\sigma[v]$.

A pair $(\sigma,\tau)$ of positional strategies turns $G$ into a finite
Markov chain $M_{\sigma,\tau}$: from $v\in\Vmax$ the token moves to
$\sigma(v)$ with probability $1$, from $v\in\Vmin$ to $\tau(v)$ with
probability $1$, and from $v\in\Vavg$ to each of $v^{(0)},v^{(1)}$ with
probability $\tfrac12$ each. Let
\[
  v_{\sigma,\tau}(u)\;:=\;\Pr\nolimits^{M_{\sigma,\tau}}_{u}\bigl[\text{the token reaches } t_1\bigr].
\]
The \emph{value} of $G$ at $u$ is
$\val(u):=\max_\sigma\min_\tau v_{\sigma,\tau}(u)$, the maximum and
minimum ranging over positional strategies. For a fixed Max strategy
$\sigma$ we write
\[
  \val_\sigma\;:=\;\min_\tau v_{\sigma,\tau}
\]
for the pointwise minimum over positional Min strategies --- the value
Max can guarantee by committing to $\sigma$. By \Cref{thm:determinacy}
applied to the one-player game obtained by freezing $\sigma$, the minimum
is attained simultaneously at all vertices by a single positional $\tau$,
so $\val_\sigma$ is itself of the form $v_{\sigma,\tau}$.
\end{definition}

\begin{problem}[\textsc{SSG-Value}]\label{prob:main}
Given an SSG $G$ and a start vertex $v_0$, decide whether
$\val(v_0)\ge\tfrac12$.
\end{problem}

The only external fact we import is the following classical determinacy
statement; everything else in this document is proved from scratch.

\begin{theorem}[Positional determinacy; Shapley, Condon]\label{thm:determinacy}
For every SSG and every vertex $u$,
\[
  \max_\sigma\min_\tau v_{\sigma,\tau}(u)
  \;=\;\min_\tau\max_\sigma v_{\sigma,\tau}(u),
\]
the extrema ranging over positional strategies, and this common value
equals the value of the game when both players are additionally allowed
arbitrary history-dependent randomised strategies. Moreover the
extrema are attained simultaneously at all vertices by a single pair
$(\sigma^\ast,\tau^\ast)$ of positional strategies.
\end{theorem}

\subsection{What is proved here, and what is not}

\Cref{prob:main} is not resolved below. What is established is a complete
and self-contained chain of results that isolates the obstruction
exactly. The logical skeleton is:

\begin{enumerate}[label=(\roman*),leftmargin=2.2em]
\item \Cref{sec:operators,sec:stopping}: the Shapley operator, its
  monotonicity, and --- for \emph{stopping} games --- a proof that its
  $N$-th iterate, and that of each one-player restriction $T_\sigma$, is a
  strict contraction, so that each has a unique fixed point
  (\Cref{thm:contraction}); determinacy and the identification of those
  fixed points with the values are then proved rather than imported
  (\Cref{thm:stopping-determinacy}). The operator $T$ itself is in general
  not a contraction.
\item \Cref{sec:transform}: an explicit, fully quantitative reduction of
  \Cref{prob:main} on an arbitrary SSG to the same problem on a stopping
  SSG of size $O(N(|\Vavg|+\log N))$ (\Cref{thm:stopping-transform}). All constants
  are computed.
\item \Cref{sec:special}: a polynomial special case
  (\Cref{thm:few-avg}) --- arbitrary SSGs with $O(\log N)$ average
  vertices --- obtained from a positional escape estimate that needs no
  stopping hypothesis.
\item \Cref{sec:structure}: the structure of the optimal set.
  \Cref{thm:opt-subcube} identifies $\Opt$ with a subcube;
  \Cref{thm:short-path} shows every strategy is joined to $\Opt$ by an
  improving path of length at most $n$. Explicit instances
  (\Cref{ex:nonstopping-nonsubcube,ex:pi-fails}) show that the stopping
  hypothesis is not removable.
\item \Cref{sec:refutations}: four refutations, each with an explicit
  instance verified in exact arithmetic, that eliminate specific
  strategies of proof.
\item \Cref{sec:selection}: a proof that the problem is polynomial-time
  equivalent to comparing two vertex values, and that $n$ non-adaptive
  comparisons determine an optimal strategy
  (\Cref{thm:compare-equivalence}).
\item \Cref{sec:randomfacet}: the strongest positive algorithmic result
  proved here --- a Las~Vegas algorithm of expected running time
  $e^{2\sqrt n}\,\mathrm{poly}(N)$ (\Cref{thm:subexp}), subexponential but
  not polynomial.
\item \Cref{sec:gap}: the exact statement of what is missing, together
  with a proof that supplying it is equivalent to resolving
  \Cref{prob:main}.
\end{enumerate}

\section{The Shapley operator}\label{sec:operators}

Throughout, vectors are indexed by $V$ and compared pointwise: $x\le y$
means $x(u)\le y(u)$ for all $u$. Let
\[
  \mathcal{X}\;:=\;\bigl\{x\in[0,1]^V \;:\; x(t_1)=1,\; x(t_0)=0\bigr\}.
\]

\begin{definition}\label{def:shapley}
The \emph{Shapley operator} $T\colon\mathcal{X}\to\mathcal{X}$ is defined
by $(Tx)(t_1)=1$, $(Tx)(t_0)=0$ and, for a non-sink $v$,
\[
  (Tx)(v)\;=\;
  \begin{cases}
    \max\bigl(x(v^{(0)}),\,x(v^{(1)})\bigr), & v\in\Vmax,\\[2pt]
    \min\bigl(x(v^{(0)}),\,x(v^{(1)})\bigr), & v\in\Vmin,\\[2pt]
    \tfrac12\bigl(x(v^{(0)})+x(v^{(1)})\bigr), & v\in\Vavg.
  \end{cases}
\]
For a Max strategy $\sigma$, the operator $T_\sigma$ is defined by the
same formulas except that the $\Vmax$ line is replaced by
$(T_\sigma x)(v)=x(\sigma(v))$. For a pair $(\sigma,\tau)$, the operator
$T_{\sigma,\tau}$ additionally replaces the $\Vmin$ line by
$(T_{\sigma,\tau}x)(v)=x(\tau(v))$; it is affine.
\end{definition}

\begin{lemma}[Monotone, nonexpansive]\label{lem:monotone}
$T$, every $T_\sigma$ and every $T_{\sigma,\tau}$ map $\mathcal{X}$ into
$\mathcal{X}$, are monotone ($x\le y\Rightarrow Tx\le Ty$), and are
nonexpansive in the supremum norm:
$\lVert Tx-Ty\rVert_\infty\le\lVert x-y\rVert_\infty$.
\end{lemma}

\begin{proof}
Each coordinate of $Tx$ is a maximum, a minimum, or a convex combination
of coordinates of $x$, and each of those three operations preserves
$[0,1]$, is monotone, and is nonexpansive: for a maximum,
$|\max(a_1,a_2)-\max(b_1,b_2)|\le\max_i|a_i-b_i|$, and symmetrically for
a minimum; for a convex combination,
$\bigl|\tfrac12(a_1+a_2)-\tfrac12(b_1+b_2)\bigr|\le\max_i|a_i-b_i|$. The
sink coordinates are pinned and contribute $0$ to the difference.
\end{proof}

\begin{lemma}[Profile values are least fixed points]\label{lem:lfp}
For every positional pair $(\sigma,\tau)$, the vector $v_{\sigma,\tau}$
is the least fixed point of $T_{\sigma,\tau}$ in $\mathcal{X}$;
equivalently, $v_{\sigma,\tau}=\min\{x\in\mathcal{X}:T_{\sigma,\tau}x\le x\}$.
\end{lemma}

\begin{proof}
Reachability probabilities in a finite Markov chain satisfy the
corresponding linear system, so $v_{\sigma,\tau}$ is a fixed point of
$T_{\sigma,\tau}$. Let $x\in\mathcal{X}$ satisfy
$T_{\sigma,\tau}x\le x$. Write $p_k(u)$ for the probability of reaching
$t_1$ from $u$ within $k$ steps; then $p_0=\bbone_{\{t_1\}}\le x$, and by
induction using monotonicity
$p_{k+1}=T_{\sigma,\tau}\,p_k\le T_{\sigma,\tau}x\le x$. Since
$p_k\uparrow v_{\sigma,\tau}$ we get $v_{\sigma,\tau}\le x$. As
$v_{\sigma,\tau}$ is itself a fixed point, it is the least one and also
the least prefixed point; this is the Knaster--Tarski characterisation
made explicit.
\end{proof}

\begin{remark}\label{rem:no-postfixed}
The dual statement fails: $x\le T_{\sigma,\tau}x$ does \emph{not} imply
$x\le v_{\sigma,\tau}$. Taking a single average vertex whose two edges
both return to itself, and $x\equiv 1$, gives $T_{\sigma,\tau}x=x$ while
$v_{\sigma,\tau}=0$. This is exactly the phenomenon that the stopping
hypothesis of the next section removes, and it is why several arguments
below are stated for stopping games. \Cref{ex:pi-fails} shows the same
failure doing real damage to an algorithm.
\end{remark}

\begin{remark}\label{rem:lfp-scope}
We deliberately do \emph{not} assert here that the value $\val$ is the
least fixed point of $T$, nor that $\val_\sigma$ is the least fixed point
of $T_\sigma$. Those statements are true, but proving them requires an
exchange of a limit with an infimum that does not follow from the
argument above. For stopping games we obtain the corresponding
identifications in \Cref{thm:stopping-determinacy} by a different route,
from uniqueness rather than from minimality; and for general games we use
only \Cref{thm:determinacy} and \Cref{lem:lfp}, never a least-fixed-point
description of $\val$.
\end{remark}

\section{Stopping games}\label{sec:stopping}

\begin{definition}\label{def:stopping}
An SSG $G$ is \emph{stopping} if for every positional pair
$(\sigma,\tau)$ and every vertex $u$, the chain $M_{\sigma,\tau}$ reaches
$\{t_0,t_1\}$ from $u$ with probability $1$.
\end{definition}

The absorption estimate below must hold not merely for positional pairs
but for arbitrary, possibly history-dependent and randomised, play: it is
applied inside a finite-horizon game whose optimal strategies are
time-dependent. We therefore prove it by an attractor argument, which is
insensitive to how the players choose.

\begin{definition}[Sink attractor]\label{def:attractor}
Put $A_0:=\{t_0,t_1\}$ and, for $j\ge0$,
\[
  A_{j+1}\;:=\;A_j\;\cup\;\bigl\{v\in\Vavg:\{v^{(0)},v^{(1)}\}\cap A_j\ne\emptyset\bigr\}
  \;\cup\;\bigl\{v\in\Vmax\cup\Vmin:\{v^{(0)},v^{(1)}\}\subseteq A_j\bigr\}.
\]
The sets $A_j$ increase, so they stabilise at $A_\infty:=A_{N}$. For
$u\in A_\infty$ write $\ell(u):=\min\{j:u\in A_j\}$.
\end{definition}

\begin{lemma}[Uniform absorption]\label{lem:absorption}
Write $a:=|\Vavg|$ and let $G$ be stopping. Then $A_\infty=V$, and for
every vertex $u$ and \emph{every} pair of strategies --- positional or
history-dependent, deterministic or randomised --- the probability of
reaching a sink from $u$ within $N$ steps is at least $2^{-a}$.
Consequently the absorption time $T_{\mathrm{abs}}$ satisfies
$\Pr_u[T_{\mathrm{abs}}>jN]\le(1-2^{-a})^j$ and
$\mathbb{E}_u[T_{\mathrm{abs}}]\le N2^{a}$.
\end{lemma}

\begin{proof}
\emph{$A_\infty=V$.} Suppose not, and set $S:=V\setminus A_\infty$, which
is nonempty and contains no sink. If $v\in S\cap\Vavg$ had a successor in
$A_\infty=A_N$ then $v$ would lie in $A_{N+1}=A_\infty$; so both
successors of such a $v$ lie in $S$. If $v\in S\cap(\Vmax\cup\Vmin)$ had
both successors in $A_\infty$ then again $v\in A_\infty$; so at least one
successor of such a $v$ lies in $S$. Choose a positional pair that at
every controlled vertex of $S$ selects a successor inside $S$ (and
arbitrarily elsewhere). Under that pair the token started in $S$ never
leaves $S$, hence never reaches a sink --- contradicting the assumption
that $G$ is stopping.

\emph{The estimate.} For a vertex $u$ let
$A(u):=\bigl|\{w\in\Vavg:\ell(w)\le\ell(u)\}\bigr|\le a$. We show by
induction on $\ell(u)$ that from $u$, under every strategy pair, a sink is
reached within $\ell(u)\le N$ steps with probability at least
$2^{-A(u)}$.

For $\ell(u)=0$ the vertex is already a sink. Let $\ell(u)=j+1$. If
$u\in\Vmax\cup\Vmin$ then \emph{both} successors lie in $A_j$, so whatever
the player does --- at any time, after any history, with any
randomisation --- the token moves in one step to some $w$ with
$\ell(w)\le j$; by induction it then reaches a sink within $j$ further
steps with probability at least $2^{-A(w)}\ge 2^{-A(u)}$, the last
inequality because $\ell(w)\le\ell(u)$ gives $A(w)\le A(u)$. If
$u\in\Vavg$ then at least one successor $w$ has $\ell(w)\le j$ and is
taken with probability $\tfrac12$, so the probability is at least
$\tfrac12\cdot2^{-A(w)}$; here $u$ itself is an average vertex with
$\ell(u)>\ell(w)$, so $u$ is counted in $A(u)$ but not in $A(w)$, giving
$A(u)\ge A(w)+1$ and hence $\tfrac12\cdot2^{-A(w)}\ge2^{-A(u)}$.

Since $A(u)\le a$ for every $u$, the probability of reaching a sink within
$\ell(u)\le N$ steps is at least $2^{-a}$, uniformly over all strategies.

Applying this to consecutive blocks of $N$ steps and using the tower
property gives $\Pr_u[T_{\mathrm{abs}}>jN]\le(1-2^{-a})^j$, whence
$\mathbb{E}[T_{\mathrm{abs}}]\le N\sum_{j\ge0}(1-2^{-a})^j=N2^{a}$.
\end{proof}

\begin{remark}
Two features of this bound are used later. First, it is
\emph{strategy-independent}, which is what allows it to be invoked inside
the finite-horizon game of \Cref{thm:contraction}, where optimal play is
time-dependent. Second, the exponent is $a=|\Vavg|$ rather than $N$: only
average vertices branch, and a level-decreasing descent meets each at
most once. The second point is what makes the constants in
\Cref{thm:stopping-transform,thm:few-avg} depend on the number of average
vertices rather than on the size of the game.
\end{remark}

\begin{theorem}[Contraction and uniqueness]\label{thm:contraction}
Let $G$ be stopping and $a=|\Vavg|$. Then for all $x,y\in\mathcal{X}$,
\[
  \lVert T^{N}x-T^{N}y\rVert_\infty\;\le\;\bigl(1-2^{-a}\bigr)\,\lVert x-y\rVert_\infty ,
\]
and the same holds for $T_\sigma$ and for $T_{\sigma,\tau}$. Consequently
each of these operators has a \emph{unique} fixed point in
$\mathcal{X}$, and every orbit converges to it. For $T_{\sigma,\tau}$
that fixed point is $v_{\sigma,\tau}$; the identification of the fixed
points of $T$ and of $T_\sigma$ with the game values is
\Cref{thm:stopping-determinacy} below, which uses only the uniqueness
established here.
\end{theorem}

\begin{proof}
Unfolding the definition $N$ times, backward induction shows that
$(T^{N}x)(u)$ is the value of the $N$-stage game started at $u$ in which
a play absorbed at $t_1$ (resp.\ $t_0$) is paid $1$ (resp.\ $0$) and a
play not absorbed within $N$ steps is paid $x(\cdot)$ at its current
vertex. In that finite-horizon game both players may use time-dependent
strategies, and backward induction produces optimal ones; this is
essential, since the value over \emph{positional} strategies alone can be
strictly smaller.

Let $\Sigma_N$ denote the set of all (history-dependent, randomised)
strategy pairs in the $N$-stage game, and for $\pi\in\Sigma_N$ write
$K_\pi(u,\cdot)$ for the law of the position after $N$ steps. The
$x$-payoff and the $y$-payoff of a fixed $\pi$ differ by
\[
  \Bigl|\sum_{w\ \mathrm{non\text{-}sink}} K_\pi(u,w)\bigl(x(w)-y(w)\bigr)\Bigr|
  \;\le\;\Pr\nolimits_{\pi}\bigl[T_{\mathrm{abs}}>N\bigr]\cdot\lVert x-y\rVert_\infty
  \;\le\;\bigl(1-2^{-a}\bigr)\lVert x-y\rVert_\infty ,
\]
where the last inequality is \Cref{lem:absorption}, which holds uniformly
over $\pi\in\Sigma_N$ precisely because its proof is
strategy-independent. Two zero-sum games on the same strategy sets whose
payoff functions differ pointwise by at most $\varepsilon$ have values
differing by at most $\varepsilon$; applying this with
$\varepsilon=(1-2^{-a})\lVert x-y\rVert_\infty$ gives the displayed
estimate. The arguments for $T_\sigma$ and $T_{\sigma,\tau}$ are
identical with one or both players' choices frozen.

A map whose $N$-th iterate is a strict contraction on the complete metric
space $\mathcal{X}$ has a unique fixed point, to which all orbits
converge. For $T_{\sigma,\tau}$ that fixed point is $v_{\sigma,\tau}$ by
\Cref{lem:lfp}. The identification of the fixed points of $T_\sigma$ and
$T$ with $\val_\sigma$ and $\val$ is \Cref{thm:stopping-determinacy}
below, which is proved from the uniqueness just established and does not
presuppose it.
\end{proof}

For stopping games we do not need to import \Cref{thm:determinacy}: both
determinacy and the identification of the fixed points with the values
follow from \Cref{thm:contraction} by a verification argument.

\begin{theorem}[Determinacy for stopping games]\label{thm:stopping-determinacy}
Let $G$ be stopping, let $w^\ast$ be the unique fixed point of $T$ and,
for a Max strategy $\sigma$, let $y_\sigma$ be the unique fixed point of
$T_\sigma$ (both exist by \Cref{thm:contraction}). Let $\sigma^\ast$
choose at each $v\in\Vmax$ a successor attaining
$\max\bigl(w^\ast(v^{(0)}),w^\ast(v^{(1)})\bigr)$ and let $\tau^\ast$
choose at each $v\in\Vmin$ a successor attaining the corresponding
minimum. Then
\[
  w^\ast\;=\;\max_\sigma\min_\tau v_{\sigma,\tau}\;=\;\min_\tau\max_\sigma v_{\sigma,\tau}\;=\;\val ,
\]
the extrema being over positional strategies and attained at
$(\sigma^\ast,\tau^\ast)$; and $y_\sigma=\min_\tau v_{\sigma,\tau}=\val_\sigma$
for every $\sigma$.
\end{theorem}

\begin{proof}
By the choice of $\sigma^\ast,\tau^\ast$ we have
$T_{\sigma^\ast,\tau^\ast}w^\ast=Tw^\ast=w^\ast$, so $w^\ast$ is a fixed
point of $T_{\sigma^\ast,\tau^\ast}$, whence
$w^\ast=v_{\sigma^\ast,\tau^\ast}$ by \Cref{thm:contraction}.

Fix any positional $\tau$. At $v\in\Vmin$ we have
$(T_{\sigma^\ast,\tau}w^\ast)(v)=w^\ast(\tau(v))\ge\min\bigl(w^\ast(v^{(0)}),w^\ast(v^{(1)})\bigr)=w^\ast(v)$,
and at every other vertex $T_{\sigma^\ast,\tau}$ agrees with
$T_{\sigma^\ast,\tau^\ast}$, so $T_{\sigma^\ast,\tau}w^\ast\ge w^\ast$.
By monotonicity the iterates $T_{\sigma^\ast,\tau}^{k}w^\ast$ are
nondecreasing, and by \Cref{thm:contraction} they converge to
$v_{\sigma^\ast,\tau}$; hence $v_{\sigma^\ast,\tau}\ge w^\ast$. As $\tau$
was arbitrary, $\min_\tau v_{\sigma^\ast,\tau}\ge w^\ast$ and therefore
$\max_\sigma\min_\tau v_{\sigma,\tau}\ge w^\ast$. Symmetrically, for
every $\sigma$ we get $T_{\sigma,\tau^\ast}w^\ast\le w^\ast$ and hence
$v_{\sigma,\tau^\ast}\le w^\ast$, so
$\min_\tau\max_\sigma v_{\sigma,\tau}\le w^\ast$. Since
$\max_\sigma\min_\tau\le\min_\tau\max_\sigma$ always, all three
quantities coincide with $w^\ast$, and the extrema are attained at
$(\sigma^\ast,\tau^\ast)$.

The statement for $y_\sigma$ is the same argument with Max's choices
frozen: let $\tau_\sigma$ be greedy for $y_\sigma$ at $\Vmin$; then
$T_{\sigma,\tau_\sigma}y_\sigma=T_\sigma y_\sigma=y_\sigma$ gives
$y_\sigma=v_{\sigma,\tau_\sigma}$, and $T_{\sigma,\tau}y_\sigma\ge y_\sigma$
for every $\tau$ gives $v_{\sigma,\tau}\ge y_\sigma$; hence
$y_\sigma=\min_\tau v_{\sigma,\tau}$, the minimum being attained by the
single positional strategy $\tau_\sigma$.
\end{proof}

\begin{corollary}[Two-sided comparison]\label{cor:comparison}
Let $G$ be stopping and $\sigma$ a Max strategy. If $x\in\mathcal{X}$
satisfies $x\le T_\sigma x$ then $x\le\val_\sigma$; if
$T_\sigma x\le x$ then $\val_\sigma\le x$. The same holds for $T$ and
$\val$.
\end{corollary}

\begin{proof}
By \Cref{thm:contraction} and \Cref{thm:stopping-determinacy},
$\val_\sigma$ is the unique fixed point of $T_\sigma$ and every orbit of
$T_\sigma$ converges to it. If $x\le T_\sigma x$ then by monotonicity the
sequence $T_\sigma^{k}x$ is nondecreasing, and it converges to
$\val_\sigma$; hence $x\le\val_\sigma$. If $T_\sigma x\le x$ then
$T_\sigma^{k}x$ is nonincreasing and converges to the same limit, so
$\val_\sigma\le x$. The argument for $T$ and $\val$ is identical.
\end{proof}

\Cref{cor:comparison} is exactly the statement that fails in general
SSGs, by \Cref{rem:no-postfixed}. It is the engine of every switching
argument below.

\section{The quantitative stopping transformation}\label{sec:transform}

\begin{definition}[Damping gadget]\label{def:damping}
Let $G$ be an SSG and $m\ge1$ an integer. The game $G_m$ is obtained by
replacing each edge $u\to w$ of $G$ by a fresh chain of $m$ average
vertices $g_1,\dots,g_m$, where $u$'s edge now leads to $g_1$,
\[
  g_i \longrightarrow \bigl\{w,\;g_{i+1}\bigr\}\quad (1\le i<m),
  \qquad
  g_m \longrightarrow \bigl\{w,\;t_0\bigr\},
\]
each with probability $\tfrac12$. The vertex sets $\Vmax$ and $\Vmin$,
and hence both players' strategy spaces, are unchanged.
\end{definition}

\begin{lemma}[The gadget is a damped edge]\label{lem:gadget}
In $G_m$, a token entering the chain replacing $u\to w$ arrives at $w$
with probability $1-\beta$ and at $t_0$ with probability $\beta$, where
$\beta:=2^{-m}$. The game $G_m$ is a stopping SSG with
$|V(G_m)| = N + 2m(N-2) < (2m+1)N$ vertices.
\end{lemma}

\begin{proof}
From $g_1$ the token reaches $w$ at step $i$ with probability $2^{-i}$
for $1\le i\le m$, and otherwise reaches $t_0$; summing,
$\sum_{i=1}^{m}2^{-i}=1-2^{-m}$. Every non-sink vertex of $G$ has two
outgoing edges, so exactly $2(N-2)$ chains are inserted. Stopping holds
because every chain leads to $t_0$ with positive probability at its last
vertex, so from every vertex, under every strategy pair, a sink is
reachable.
\end{proof}

\begin{lemma}[Denominator bound]\label{lem:denominator}
Let $H$ be an SSG with $r:=|\Vavg|$ average vertices and let
$(\sigma,\tau)$ be a positional pair. Every coordinate of
$v_{\sigma,\tau}$ is a rational number whose denominator divides a
positive integer bounded by $8^{r/2}$. The same bound applies to $\val$,
since by \Cref{thm:determinacy} $\val=v_{\sigma^\ast,\tau^\ast}$ for a
positional pair. In particular the bound is independent of the number of
Max and Min vertices, and for $r=0$ all values lie in $\{0,1\}$.
\end{lemma}

\begin{proof}
In the chain $M_{\sigma,\tau}$ every vertex outside
$\Vavg\cup\{t_0,t_1\}$ has a single successor. For a vertex $u$ define
$D(u)$ by following those deterministic steps from $u$ until an element
of $\Vavg\cup\{t_0,t_1\}$ is reached, where we set $D(u)=u$ if $u$
already lies in that set; if the deterministic walk instead enters a
cycle disjoint from $\Vavg\cup\{t_0,t_1\}$, the play never reaches $t_1$
and we put $D(u)=\bot$. Writing $\mathcal{V}(t_1)=1$,
$\mathcal{V}(t_0)=\mathcal{V}(\bot)=0$ and
$\mathcal{V}(w)=v_{\sigma,\tau}(w)$ for $w\in\Vavg$, the strong Markov
property gives $v_{\sigma,\tau}(u)=\mathcal{V}(D(u))$ for every $u$, and
for $u\in\Vavg$ with successors $u^{(0)},u^{(1)}$,
\[
  2\,v_{\sigma,\tau}(u)\;=\;\mathcal{V}\bigl(D(u^{(0)})\bigr)+\mathcal{V}\bigl(D(u^{(1)})\bigr).
\]
Thus the restriction of $v_{\sigma,\tau}$ to $\Vavg$ satisfies a linear
system $Ax=c$ of size at most $r$, where $c$ has entries in
$\{0,1,2\}$ and $A$ has integer entries: the row of $u$ carries
$2-\mu_u(u)$ on the diagonal and $-\mu_u(w)$ in column $w\ne u$, with
$\mu_u(w):=|\{i: D(u^{(i)})=w\}|\in\{0,1,2\}$ and
$\sum_{w}\mu_u(w)\le2$. Every row therefore has Euclidean norm at most
$\sqrt{8}$, the extreme case being $\mu_u(w)=2$ for a single $w\ne u$.

Restrict the system to the set $U\subseteq\Vavg$ of average vertices from
which $t_1$ is reachable in $M_{\sigma,\tau}$; off $U$ the value is $0$,
so those coordinates may be deleted, and the resulting matrix $A_U$ is
nonsingular (from every vertex of $U$ the chain leaves $U$ with positive
probability, so the corresponding substochastic matrix has spectral
radius below $1$). Hadamard's inequality gives
$|\det A_U|\le 8^{|U|/2}\le 8^{r/2}$, and by Cramer's rule every
coordinate of $x$ is an integer divided by $\det A_U$. Since every other
value equals some $\mathcal{V}(D(u))\in\{0,1\}\cup\{x_w\}$, the whole
vector has denominators dividing $\det A_U$. If $r=0$ every value is
$\mathcal{V}(D(u))\in\{0,1\}$.
\end{proof}

\begin{theorem}[Exact reduction to a stopping game]\label{thm:stopping-transform}
Let $G$ be an SSG on $N\ge3$ vertices with $a=|\Vavg|$ average vertices
and start vertex $v_0$, put $D:=2^{\lceil 3a/2\rceil}$ --- a power of two,
and by \Cref{lem:denominator} an upper bound on the denominator of
$\val_G(v_0)$, since $8^{a/2}=2^{3a/2}\le D$ --- and choose
\[
  m \;:=\; \bigl\lceil\, \tfrac52 a + \log_2 N \,\bigr\rceil + 3 .
\]
Then $G_m$ is a stopping SSG with $O\bigl(N(a+\log N)\bigr)=O(N^2)$ vertices, computable in time
polynomial in $N$, and
\[
  \val_G(v_0)\;\ge\;\tfrac12
  \qquad\Longleftrightarrow\qquad
  \val_{G_m}(v_0)\;>\;\tfrac12-\tfrac{1}{4D}.
\]
\end{theorem}

\begin{proof}
Write $\beta=2^{-m}$. \emph{Step 1: $\val_{G_m}\le\val_G$.} Fix any pair
$(\sigma,\tau)$. Couple the two chains by letting them make identical
choices at $\Vmax\cup\Vmin\cup\Vavg$ vertices; the chain in $G_m$
additionally has, at each traversal of an edge, probability $\beta$ of
being diverted to $t_0$. Under this coupling every play of $G_m$ that
reaches $t_1$ corresponds to a play of $G$ that reaches $t_1$, so
$v^{G_m}_{\sigma,\tau}(v_0)\le v^{G}_{\sigma,\tau}(v_0)$; taking
$\max_\sigma\min_\tau$ preserves the inequality.

\emph{Step 2: $\val_G-\val_{G_m}\le\beta N2^{a}$.} Let $\sigma^\ast$ be
optimal for Max in $G$ and let $\tau$ be arbitrary. In the chain
$M_{\sigma^\ast,\tau}$ of $G$, call a vertex \emph{live} if $t_1$ is
reachable from it, and let $T_{\mathrm{res}}$ be the first time the token
reaches $t_1$ or leaves the live set. From every live vertex there is a
simple path to $t_1$, of length at most $N-1$, and by the argument of
\Cref{lem:absorption} it is followed with probability at least $2^{-a}$;
hence $\Pr[T_{\mathrm{res}}>jN]\le(1-2^{-a})^{j}$ and
$\mathbb{E}[T_{\mathrm{res}}]\le N2^{a}$. Under the coupling of Step 1, a
play of $G_m$ fails to reach $t_1$ only if the corresponding play of $G$
fails to, or if a diversion occurs before time $T_{\mathrm{res}}$. By a
union bound the latter has probability at most
$\beta\,\mathbb{E}[T_{\mathrm{res}}]\le\beta N2^{a}$. Therefore
$v^{G_m}_{\sigma^\ast,\tau}(v_0)\ge v^{G}_{\sigma^\ast,\tau}(v_0)-\beta N2^{a}$
for every $\tau$, and taking $\min_\tau$ and then the maximum over Max
strategies gives $\val_{G_m}(v_0)\ge\val_G(v_0)-\beta N2^{a}$.

\emph{Step 3: separation.} By \Cref{lem:denominator}, $\val_G(v_0)$ is a
rational with denominator at most $D$. If $\val_G(v_0)<\tfrac12$ then
$\tfrac12-\val_G(v_0)\ge\tfrac{1}{2D}$, so by Step 1
$\val_{G_m}(v_0)\le\tfrac12-\tfrac1{2D}<\tfrac12-\tfrac1{4D}$. If
$\val_G(v_0)\ge\tfrac12$ then by Step 2
$\val_{G_m}(v_0)\ge\tfrac12-\beta N2^{a}$, and the choice of $m$ gives
$\beta N 2^{a}=2^{-m}N2^{a}<\tfrac{1}{4D}$: indeed
$2^{-m}N2^{a}<2^{-2}2^{-\lceil 3a/2\rceil}$ is equivalent to
$m>\lceil\tfrac32a\rceil+a+\log_2 N+2$, and since
$\lceil\tfrac32a\rceil\le\tfrac32a+\tfrac12$ the right-hand side is at most
$\tfrac52a+\log_2N+\tfrac52$, which $m=\lceil \tfrac52a+\log_2N\rceil+3$
exceeds. Hence
$\val_{G_m}(v_0)>\tfrac12-\tfrac1{4D}$. The two cases are exclusive and
exhaustive, proving the equivalence. Finally $|V(G_m)|<(2m+1)N=O\bigl(N(a+\log N)\bigr)$ by
\Cref{lem:gadget}.
\end{proof}

\begin{remark}
\Cref{thm:stopping-transform} is what licenses us to assume that the input
game is stopping. The cost is not zero and must be tracked in two ways.
First, the threshold moves: the transformed game is compared against
$\tfrac12-\tfrac1{4D}$, not against $\tfrac12$, and forgetting this gives
wrong answers (see the proof of \Cref{thm:compare-equivalence}). Second,
the gadget inserts $\Theta(a+\log N)$ average vertices in front of every
edge, so any bound whose parameter is $|\Vavg|$ --- such as
\Cref{thm:few-avg} --- must be proved before the transformation, not
after. Every statement in
\Cref{sec:structure,sec:refutations} is therefore stated for stopping
games without loss of generality for the purposes of \Cref{prob:main}.
\end{remark}

\section{A polynomial special case}\label{sec:special}

The sharpened escape estimate has an algorithmic consequence. It is
modest, but it is a genuine positive result, it isolates $|\Vavg|$ rather
than $N$ as the parameter controlling value iteration, and --- unlike
most of what follows --- it needs no stopping hypothesis, so it does not
pay the price of \Cref{thm:stopping-transform} (which inflates $|\Vavg|$
to $\Theta(N^2)$).

\begin{lemma}[Forced descent]\label{lem:descent}
Let $G$ be an arbitrary SSG, $a=|\Vavg|$, and let $\sigma^\ast$ be an
optimal positional Max strategy (\Cref{thm:determinacy}). Put
$C_0:=\{t_1\}$ and
\[
  C_{j+1}:=C_j\cup\bigl\{v\in\Vavg:\{v^{(0)},v^{(1)}\}\cap C_j\ne\emptyset\bigr\}
  \cup\bigl\{v\in\Vmax:\sigma^\ast(v)\in C_j\bigr\}
  \cup\bigl\{v\in\Vmin:\{v^{(0)},v^{(1)}\}\subseteq C_j\bigr\},
\]
and $C_\infty:=C_{N}$. Then
\begin{enumerate}[label=(\alph*),leftmargin=2.2em]
\item $C_\infty=\{v:\val(v)>0\}$;
\item for every $v\in C_\infty$ and \emph{every} Min strategy $\pi$ ---
  positional or history-dependent, deterministic or randomised --- the
  play from $v$ under $(\sigma^\ast,\pi)$ reaches $t_1$ within $N$ steps
  with probability at least $2^{-a}$.
\end{enumerate}
\end{lemma}

\begin{proof}
(b) Write $\ell(v)$ for the least $j$ with $v\in C_j$, and
$A(v):=|\{w\in\Vavg:\ell(w)\le\ell(v)\}|\le a$. We show by induction on
$\ell(v)$ that $t_1$ is reached within $\ell(v)\le N$ steps with
probability at least $2^{-A(v)}$, whatever $\pi$ does. For $\ell(v)=0$
the vertex is $t_1$. Let $\ell(v)=j+1$. If $v\in\Vmax$ then
$\sigma^\ast(v)\in C_j$, and Max moves there with probability $1$. If
$v\in\Vmin$ then \emph{both} successors lie in $C_j$, so whatever $\pi$
does --- at any time, after any history --- the token moves to a vertex
of level at most $j$. If $v\in\Vavg$ then at least one successor lies in
$C_j$ and is taken with probability $\tfrac12$; here $v$ is itself an
average vertex of level $j+1$, so it is counted in $A(v)$ but not in
$A(w)$ for the successor $w$, giving $A(v)\ge A(w)+1$ and hence
$\tfrac12\cdot2^{-A(w)}\ge2^{-A(v)}$. In every case the token descends to
a strictly smaller level, so it reaches $t_1$ within $\ell(v)$ steps.

(a) ($\supseteq$) By (b), $\val(v)\ge\val_{\sigma^\ast}(v)\ge2^{-a}>0$ for
$v\in C_\infty$. ($\subseteq$) Let $S:=V\setminus C_\infty$; then
$t_1\notin S$, every $v\in S\cap\Vmax$ has $\sigma^\ast(v)\in S$, every
$v\in S\cap\Vavg$ has both successors in $S$, and every $v\in S\cap\Vmin$
has at least one successor in $S$ --- in each case because the opposite
would place $v$ in $C_\infty$. Let Min play, inside $S$, a successor in
$S$. Then from any $v\in S$ the token never leaves $S$, hence never
reaches $t_1$, so $\val(v)\le\val_{\sigma^\ast}(v)=0$.
\end{proof}

\begin{theorem}\label{thm:few-avg}
Every SSG on $N$ vertices with $a=|\Vavg|$ average vertices, stopping or
not, can be decided --- and its exact value vector computed --- in time
$\mathrm{poly}(N)\cdot 2^{a}$. In particular SSGs with $a=O(\log N)$
average vertices are decidable in polynomial time, and SSGs with no
average vertices at all in time $\mathrm{poly}(N)$.
\end{theorem}

\begin{proof}
Fix an optimal positional $\sigma^\ast$ and let $F:=T_{\sigma^\ast}$, the
operator with Max's choices frozen. Put
$u_j:=F^{jN}\bigl(\bbone_{\{t_1\}}\bigr)$, so that $u_j(v)$ is the
infimum, over all Min strategies, of the probability of reaching $t_1$
from $v$ within $jN$ steps; the infimum is attained by backward
induction, and Min may use time-dependent play. Since $T\ge F$ pointwise
and both are monotone,
\[
  \bbone_{\{t_1\}}\ \le\ u_j\ \le\ T^{jN}\bigl(\bbone_{\{t_1\}}\bigr)\ \le\ \val,
\]
the last inequality because reaching $t_1$ within $jN$ steps implies
reaching it at all. So it suffices to bound $\val-u_j$.

Off $C_\infty$ we have $\val=0$ by \Cref{lem:descent}(a) and hence
$u_j=0$ there too. Put $M_j:=\max_{v\in C_\infty}(\val-u_j)(v)$, so
$M_0\le1$. Because $\sigma^\ast$ is optimal, $\val=\val_{\sigma^\ast}$ is
the value of the Markov decision process with Max frozen and therefore
satisfies $F\val=\val$, hence $F^{N}\val=\val$. Now $F^{N}(y)(v)$ equals
$\inf_\pi\bigl[\Pr^{\pi}_v[T\le N]+\mathbb{E}^{\pi}_v[y(X_N)\mathbf1\{T>N\}]\bigr]$,
where $T$ is the hitting time of $t_1$; since
$\inf_\pi A(\pi)-\inf_\pi B(\pi)\le\sup_\pi(A(\pi)-B(\pi))$,
\[
  (\val-u_{j+1})(v)\;=\;\bigl(F^{N}\val-F^{N}u_j\bigr)(v)
  \;\le\;\sup_\pi\ \mathbb{E}^{\pi}_v\bigl[(\val-u_j)(X_N)\,\mathbf1\{T>N\}\bigr].
\]
The integrand vanishes off $C_\infty$ and is at most $M_j$ on it, so for
$v\in C_\infty$ the right-hand side is at most
$M_j\cdot\sup_\pi\Pr^{\pi}_v[T>N]\le M_j(1-2^{-a})$ by
\Cref{lem:descent}(b), the supremum being over \emph{all} Min strategies
--- which is exactly what that lemma provides. Hence
$M_{j+1}\le(1-2^{-a})M_j$ and $\val-u_j\le(1-2^{-a})^{j}\le e^{-j2^{-a}}$
everywhere.

Put $D:=2^{\lceil 3a/2\rceil}$, which bounds the denominators of $\val$ by
\Cref{lem:denominator} and \Cref{thm:determinacy}. Choosing
$j:=\bigl\lceil 2^{a}(2+\lceil\tfrac32a\rceil)\ln2\bigr\rceil+1$ makes the
error below $2^{-2-\lceil 3a/2\rceil}=1/(4D)$. If $\val(v_0)<\tfrac12$
then, $\val(v_0)$ being a rational of denominator at most $D$, we have
$\val(v_0)\le\tfrac12-\tfrac1{2D}$ and hence
$T^{jN}(\bbone_{\{t_1\}})(v_0)\le\tfrac12-\tfrac1{2D}$; if
$\val(v_0)\ge\tfrac12$ then
$T^{jN}(\bbone_{\{t_1\}})(v_0)>\tfrac12-\tfrac1{4D}$. So the threshold
$\tfrac12-\tfrac1{4D}$ decides. The exact value vector is recovered by
rounding each coordinate to the unique rational of denominator at most
$D$ within distance $1/(2D^{2})$, which needs accuracy $1/(2D^2)$ ---
double $\lceil\tfrac32a\rceil$ in $j$ --- followed by the
continued-fraction algorithm.

For the running time, iterate $T$ in binary fixed-point arithmetic with
$B$ bits after the point, rounding every intermediate value down.
Rounding down preserves $T^{k}(\bbone_{\{t_1\}})\le\val$ and adds at most
$2^{-B}$ per application of $T$, hence at most $jN2^{-B}$ in all;
$B=\Theta(a+\log N)$ makes this negligible against $1/(4D)$. The number of
applications of $T$ is $jN=O(2^{a}aN)$, each costing $O(N)$ additions,
comparisons and halvings of $B$-bit numbers, for a total of
$\mathrm{poly}(N)\cdot2^{a}$. Note that the algorithm itself iterates $T$
and never needs $\sigma^\ast$, which appears only in the analysis.
\end{proof}

\begin{remark}
\Cref{thm:few-avg} should be read together with \Cref{thm:vi-lower},
which exhibits games where value iteration provably needs
$2^{\Omega(N)}$ steps; there \emph{every} non-sink vertex is an average
vertex, so $a=N-2$ and the two results are consistent and jointly tight
in their dependence on $a$. Note also that routing this theorem through
\Cref{thm:stopping-transform} would destroy it: the damping gadget
inserts $\Theta(N)$ fresh average vertices in front of every edge, raising
$a$ to $\Theta(N^{2})$. It is precisely because \Cref{lem:descent} needs no
stopping hypothesis that the theorem applies to arbitrary games.
\end{remark}

A third tractable regime is delimited not by how many average vertices
there are, but by where they sit.

\begin{definition}[Average-jump digraph]\label{def:jump}
For an SSG $G$, let $D(G)$ be the digraph on vertex set $\Vavg$ with an
edge $u\to z$ whenever some path $u=p_0\to p_1\to\cdots\to p_k=z$ in $G$
has $k\ge1$, $z\in\Vavg$ and $p_1,\dots,p_{k-1}\notin\Vavg$.
\end{definition}

\begin{lemma}\label{lem:jump-acyclic}
If no cycle of $G$ contains an average vertex, then $D(G)$ is acyclic.
\end{lemma}

\begin{proof}
A cycle $u_1\to\cdots\to u_k\to u_1$ of $D(G)$ expands, edge by edge,
into a closed walk of $G$ through $u_1\in\Vavg$. A shortest closed walk
of positive length from $u_1$ to itself repeats no vertex --- splitting
at a repetition, including a repetition of $u_1$, would produce a shorter
positive-length closed walk at $u_1$ --- hence is a cycle, and it
contains the average vertex $u_1$, a contradiction.
\end{proof}

\begin{lemma}[Deterministic games with terminal payoffs]\label{lem:det-game}
Let $H$ be a game whose non-terminal vertices are partitioned into Max
and Min vertices, each of out-degree two, with a terminal set $Z$
carrying payoffs $c:Z\to[0,1]\cap\mathbb{Q}$, where a play never reaching
$Z$ pays Max $0$. For $\theta>0$ write
$X_\theta:=\{z\in Z:c(z)\ge\theta\}$ and let $\mathrm{Attr}(X)$ be the set
of vertices from which Max can force a visit to $X$. Then the value of
$H$ at $v$ equals
$\max\bigl(\{0\}\cup\{\theta\in c(Z):v\in\mathrm{Attr}(X_\theta)\}\bigr)$,
computable in time $O(|Z|\cdot|H|)$.
\end{lemma}

\begin{proof}
If $v\in\mathrm{Attr}(X_\theta)$ then Max forces a visit to $X_\theta$
within $|H|$ steps and is paid at least $\theta$. If
$v\notin\mathrm{Attr}(X_\theta)$ then the complement of
$\mathrm{Attr}(X_\theta)$ is a trap for Min: every Max vertex outside it
has both successors outside, and every Min vertex outside it has one, so
Min has a positional strategy keeping the play out of $X_\theta$ forever.
The play then either never reaches $Z$, paying $0<\theta$, or reaches
some $z$ with $c(z)<\theta$; either way the value is below $\theta$. The
value is therefore the largest $\theta$ in the finite set $c(Z)$ whose
attractor contains $v$, or $0$ if none does. Each attractor is one
backward sweep.
\end{proof}

\begin{theorem}[Average-acyclic games are easy]\label{thm:avg-acyclic}
Let $G$ be an SSG in which no cycle contains an average vertex. Then
$\val$ is computable exactly in polynomial time.
\end{theorem}

\begin{proof}
By \Cref{lem:jump-acyclic}, $D(G)$ is acyclic; process $\Vavg$ in reverse
topological order of $D(G)$. Suppose $u\in\Vavg$ is next and every
$D(G)$-successor of $u$ already carries its correct value. Form $H_u$
from $G$ by declaring
$Z:=\{t_0,t_1\}\cup\{\text{already-solved average vertices}\}$ terminal,
with payoffs $1$, $0$ and the computed values. Every average vertex
reachable from a successor of $u$ by a path avoiding $\Vavg$ in its
interior is a $D(G)$-successor of $u$, hence lies in $Z$; so the part of
$H_u$ \emph{reachable from $u$'s successors} is genuinely deterministic.
Elsewhere $H_u$ may still contain unsolved average vertices, so
\Cref{lem:det-game} is applied to that reachable part --- legitimate,
because whether a vertex lies in an attractor depends only on the
vertices reachable from it.

We claim the value of $H_u$ at such a vertex $v$ equals $\val(v)$. For
``$\ge$'': by \Cref{lem:det-game} Max can force a visit to $X_\theta$ and
then continue with a strategy optimal in $G$ from the terminal $z$
reached, obtaining at least $c(z)=\val(z)\ge\theta$. For ``$\le$'': if
$v\notin\mathrm{Attr}(X_\theta)$, Min keeps the play out of $X_\theta$ and,
on reaching a terminal $z$, switches to a strategy optimal in $G$ from
$z$; the play then either never reaches $t_1$, worth $0$, or is held to
$\val(z)<\theta$. Both use only \Cref{thm:determinacy} and the strong
Markov property.

Finally $\val(u)=\tfrac12\bigl(\val(u^{(0)})+\val(u^{(1)})\bigr)$, with
both terms now known, which advances the induction. After all of $\Vavg$
is processed, one further application of \Cref{lem:det-game} gives the
values at the remaining vertices. There are at most $|\Vavg|+1$ rounds,
each costing $O(N^{2})$, so the total is polynomial.
\end{proof}

The dual restriction --- forbidding \emph{controlled} vertices on cycles
rather than average ones --- is also tractable, for a completely
different reason.

\begin{theorem}[Player-free cycles]\label{thm:player-free}
Let $G$ be an SSG in which no cycle contains a vertex of
$\Vmax\cup\Vmin$. Then $\val$ is computable exactly in polynomial time.
\end{theorem}

\begin{proof}
Condense $G$ into strongly connected components and process them in
reverse topological order, so that when a component is treated, the
values at all vertices outside it that it can reach are already known.

A component that is a single vertex without a self-loop has both
successors in later-processed components, so its value is obtained
directly as the maximum, minimum or average of two known numbers.

Any other component $C$ --- non-trivial, or a single vertex with a
self-loop --- contains a cycle through each of its vertices, so by
hypothesis $C\subseteq\Vavg$. There are no choices inside $C$: the
restriction of the play to $C$ is a Markov chain with exits to known
values. Let $U\subseteq C$ be the set of vertices from which $t_1$ is
reachable in $G$; for $u\in C\setminus U$ we have $\val(u)=0$. For
$u\in U$ the value satisfies
\[
  \val(u)\;=\;\tfrac12\!\!\sum_{w\in\mathrm{succ}(u)}\!\!\val(w)
  \;=\;\sum_{w\in U}P(u,w)\val(w)\;+\;b(u),
\]
where $P$ is the substochastic matrix of the moves that stay inside $U$
and $b(u)$ collects the already-known values of the exits. From every
$u\in U$ some path leads to $t_1$, and since $t_1\notin C$ that path
leaves $U$; so the chain leaves $U$ with positive probability within
$|C|$ steps, the spectral radius of $P$ is below $1$, and $I-P$ is
invertible. The values on $U$ are therefore the unique solution of a
linear system, obtained by Gaussian elimination.

Each of the at most $N$ components costs $O(N^{3})$ exact rational
operations on numbers of polynomial bit-size (\Cref{lem:denominator}), so
the whole computation is polynomial.
\end{proof}

\begin{remark}[The two tractable cycle regimes]\label{rem:two-regimes}
\Cref{thm:avg-acyclic,thm:player-free} are complementary:
the first forbids average vertices on cycles, the second forbids
controlled ones. Each is easy for its own reason --- the first because
what remains is a finite sequence of deterministic games with terminal
payoffs, the second because what remains is a finite sequence of Markov
chains --- and together they say that \Cref{prob:main} is easy as soon as
the cyclic structure of the game is \emph{pure}. Difficulty requires
cycles on which randomness and choice act together. The point is sharpened
by \Cref{thm:vi-lower}: its family has only average vertices, so every
cycle in it is player-free and \Cref{thm:player-free} solves it in
polynomial time by one linear solve --- even though value iteration on it
provably needs $2^{\Omega(N)}$ steps. Slow value iteration is therefore
not the same thing as a hard instance. The algorithm of
\Cref{thm:player-free} was verified against exhaustive computation on
$2239$ instances with no discrepancy.

Each theorem is the degenerate case of a natural parameter --- the
minimum number of vertices whose deletion leaves no cycle through an
average vertex, respectively through a controlled vertex --- and one
expects the tractability to degrade gracefully in that parameter. We do
not develop either parameterisation here, and nothing below depends on
it.
\end{remark}

\begin{remark}
\Cref{thm:avg-acyclic} strictly extends the case $a=0$ of
\Cref{thm:few-avg}: a game with no average vertices trivially has no
cycle through one, but the converse fails --- an average-acyclic game may
have arbitrarily many average vertices, as long as none of them lies on a
cycle. The two tractable regimes are incomparable: \Cref{thm:few-avg}
allows average vertices anywhere but only $O(\log N)$ of them, while
\Cref{thm:avg-acyclic} allows any number but forbids them on cycles.
Both were verified against exhaustive computation, the algorithm of
\Cref{thm:avg-acyclic} on $3477$ average-acyclic instances with no
discrepancy.
\end{remark}

The other boundary of tractability is the number of players.

\begin{theorem}[One-player games are easy]\label{thm:one-player}
Let $G$ be an SSG with $\Vmin=\emptyset$, stopping or not. Then
$\val=\lfp(T)$, and $\val$ is the unique optimum of the linear program
\[
\begin{aligned}
  \text{minimise } &\textstyle\sum_{u\in V}x(u)\quad\text{subject to } x\in[0,1]^V,\ x(t_1)=1,\ x(t_0)=0,\\
  &x(v)\ \ge\ x(v^{(0)}),\quad x(v)\ \ge\ x(v^{(1)}) && (v\in\Vmax),\\
  &x(v)\ =\ \tfrac12\bigl(x(v^{(0)})+x(v^{(1)})\bigr) && (v\in\Vavg).
\end{aligned}
\]
Hence $\val$ is computable exactly in polynomial time. The same holds
when $\Vmax=\emptyset$ and $G$ is stopping, by \Cref{lem:duality}.
\end{theorem}

\begin{proof}
Put $q_k:=T^{k}\bigl(\bbone_{\{t_1\}}\bigr)$, so that $q_k(v)$ is the
supremum over Max strategies of the probability of reaching $t_1$ from
$v$ within $k$ steps. Suprema commute, so
\[
  \lim_k q_k(v)=\sup_k\sup_\sigma\Pr\nolimits^{\sigma}_v[\text{within }k]
  =\sup_\sigma\sup_k\Pr\nolimits^{\sigma}_v[\text{within }k]
  =\sup_\sigma\Pr\nolimits^{\sigma}_v[\text{reach }t_1]=\val(v),
\]
the last step by \Cref{thm:determinacy}. (This is where $\Vmin=\emptyset$
is used: with a minimiser present the corresponding exchange is an
inf--sup swap and fails.) Since $T$ is continuous, $T\val=\lim_kTq_k=\lim_kq_{k+1}=\val$,
so $\val$ is a fixed point. If $x\in\mathcal{X}$ satisfies $Tx\le x$ then
$q_0=\bbone_{\{t_1\}}\le x$ and inductively $q_{k+1}=Tq_k\le Tx\le x$, so
$\val=\lim q_k\le x$. Hence $\val=\lfp(T)=\min\{x\in\mathcal{X}:Tx\le x\}$,
the minimum being pointwise.

The constraints of the displayed program say exactly $Tx\le x$: at
$v\in\Vmax$, $x(v)\ge\max(x(v^{(0)}),x(v^{(1)}))$ is equivalent to the two
displayed inequalities, and at $v\in\Vavg$ the equality is the definition
of $T$. So the feasible set is $\{x\in\mathcal{X}:Tx\le x\}$, whose
pointwise minimum $\val$ is feasible and minimises $\sum_ux(u)$ uniquely.
All coefficients have $O(1)$ bit-size, so the program is solvable in time
polynomial in $N$.
\end{proof}

\begin{remark}[What one-player tractability does and does not rule out]\label{rem:owner-blind}
\Cref{thm:one-player} constrains the shape a \emph{hardness certificate}
can take. Let $\mu(G)$ be any quantity depending only on the digraph of
$G$ and on the partition of its vertices into \emph{controlled}
($\Vmax\cup\Vmin$) and \emph{average} ones --- not on which controlled
vertices belong to Max and which to Min. Then $\mu$ takes the same value
on a two-player game and on the one-player game obtained by transferring
every controlled vertex to Max, and the latter is solvable in polynomial
time. Hence a large value of such a $\mu$ can never certify that an
instance is hard: it is large on instances that are easy for reasons
$\mu$ cannot see. This is the correct reading of the observation that the
condition measures proposed for interior-point methods are blind to
ownership.

It would be a non sequitur to conclude that owner-blind parameters are
useless, and this document contains two counterexamples to that stronger
claim. The parameter $\mu(G):=2^{|\Vavg|}$ is owner-blind and exponential
on some one-player games, yet \Cref{thm:few-avg} gives a
$\mathrm{poly}(N,\mu)$ algorithm for arbitrary two-player games; and the
hypothesis of \Cref{thm:avg-acyclic} --- no cycle through an average
vertex --- is a property of the digraph and the average set alone, and
yields polynomial time for two-player games. So an owner-blind parameter
can perfectly well drive a useful algorithm. What it cannot do is witness
hardness.
\end{remark}

\begin{remark}
\Cref{thm:few-avg} should be read together with \Cref{thm:vi-lower},
which exhibits games where value iteration provably needs
$2^{\Omega(N)}$ steps; there \emph{every} non-sink vertex is an average
vertex, so $a=N-2$ and the two results are consistent and jointly tight
in their dependence on $a$. Note also that routing this theorem through
\Cref{thm:stopping-transform} would destroy it: the damping gadget
inserts $\Theta(N)$ fresh average vertices in front of every edge, raising
$a$ to $\Theta(N^{2})$. It is precisely because \Cref{lem:descent} needs no
stopping hypothesis that the theorem applies to arbitrary games.
\end{remark}

\section{The structure of the optimal set}\label{sec:structure}

Throughout this section $G$ is a \emph{stopping} SSG with value vector
$w^\ast:=\val$, and
\[
  \Opt \;:=\;\bigl\{\sigma \;:\; \val_\sigma = w^\ast\bigr\}
\]
is the set of value-optimal positional Max strategies. Identifying
$\sigma$ with its indicator bit vector in $\{0,1\}^{\Vmax}$, we may speak
of subcubes: a \emph{subcube} is a set of the form
$\prod_{v\in\Vmax}A_v$ with $\emptyset\ne A_v\subseteq\{0,1\}$.

\begin{definition}[Switchable vertices]\label{def:switchable}
A vertex $v\in\Vmax$ is \emph{strictly switchable} at $\sigma$ if
$\val_\sigma(\bar\sigma(v))>\val_\sigma(\sigma(v))$. Write $S_\sigma$ for
the set of strictly switchable vertices at $\sigma$.
\end{definition}

\begin{lemma}[Local optimality equals global optimality]\label{lem:local-global}
Let $G$ be stopping. Then $S_\sigma=\emptyset$ if and only if
$\sigma\in\Opt$.
\end{lemma}

\begin{proof}
If $S_\sigma=\emptyset$ then for every $v\in\Vmax$ we have
$\val_\sigma(v)=\val_\sigma(\sigma(v))\ge\val_\sigma(\bar\sigma(v))$,
so $\val_\sigma(v)=\max\bigl(\val_\sigma(v^{(0)}),\val_\sigma(v^{(1)})\bigr)$;
at $\Vmin$ and $\Vavg$ vertices $\val_\sigma$ already satisfies the
defining equations of $T$. Hence $T\val_\sigma=\val_\sigma$, and by
uniqueness (\Cref{thm:contraction}) $\val_\sigma=w^\ast$.

Conversely let $\sigma\in\Opt$, so $\val_\sigma=w^\ast$. Since
$w^\ast=Tw^\ast$, for $v\in\Vmax$ we have
$w^\ast(v)=\max\bigl(w^\ast(v^{(0)}),w^\ast(v^{(1)})\bigr)$, while
$\val_\sigma=T_\sigma\val_\sigma$ gives $w^\ast(v)=w^\ast(\sigma(v))$.
Hence $w^\ast(\sigma(v))\ge w^\ast(\bar\sigma(v))$ and $v\notin S_\sigma$.
\end{proof}

\begin{theorem}[The optimal set is a subcube]\label{thm:opt-subcube}
Let $G$ be a stopping SSG with value vector $w^\ast$. Then
\[
  \Opt\;=\;\Bigl\{\sigma \;:\; w^\ast(\sigma(v))=\max\bigl(w^\ast(v^{(0)}),\,w^\ast(v^{(1)})\bigr)\ \text{ for every } v\in\Vmax\Bigr\}.
\]
In particular $\Opt$ is a nonempty subcube of $\{0,1\}^{\Vmax}$; it is
therefore a sublattice, it is connected under single-coordinate flips,
and the choices of an optimal strategy at distinct vertices may be
combined freely.
\end{theorem}

\begin{proof}
($\subseteq$) Let $\sigma\in\Opt$. The second paragraph of the proof of
\Cref{lem:local-global} shows exactly that $w^\ast(\sigma(v))$ attains
the maximum, for every $v\in\Vmax$.

($\supseteq$) Suppose $\sigma$ attains the maximum at every $v\in\Vmax$.
We claim $T_\sigma w^\ast=w^\ast$. At $v\in\Vmax$,
$(T_\sigma w^\ast)(v)=w^\ast(\sigma(v))=\max\bigl(w^\ast(v^{(0)}),w^\ast(v^{(1)})\bigr)=(Tw^\ast)(v)=w^\ast(v)$
by hypothesis and $w^\ast=Tw^\ast$. At $\Vmin$ and $\Vavg$ vertices and
at the sinks, $T_\sigma$ and $T$ are given by identical formulas, so
again $(T_\sigma w^\ast)(v)=(Tw^\ast)(v)=w^\ast(v)$. Since $G$ is
stopping, $T_\sigma$ has a unique fixed point, namely $\val_\sigma$
(\Cref{thm:contraction}); hence $\val_\sigma=w^\ast$ and
$\sigma\in\Opt$. Nonemptiness is \Cref{thm:stopping-determinacy}. The description
exhibits $\Opt$ as the product over $v\in\Vmax$ of the nonempty set of
maximising choices at $v$, which is the asserted subcube.
\end{proof}

The stopping hypothesis in \Cref{thm:opt-subcube} cannot be dropped.

\begin{example}[$\Opt$ need not be a subcube without stopping]\label{ex:nonstopping-nonsubcube}
The failure already occurs for a \emph{one-player, deterministic} game,
so it has nothing to do with randomness or with the second player. Let
$V=\{v_0,v_1,v_2,v_3\}\cup\{t_0,t_1\}$ with
$\Vmax=\{v_0,v_1,v_2,v_3\}$ and $\Vmin=\Vavg=\emptyset$, and edges
\[
  v_0\to\{v_3,v_2\},\qquad
  v_1\to\{t_1,v_0\},\qquad
  v_2\to\{v_3,v_0\},\qquad
  v_3\to\{v_0,v_1\}.
\]
The game is not stopping: the strategy sending $v_0\to v_2$ and
$v_2\to v_0$ traps the token in the two-cycle $v_0\to v_2\to v_0$. Since
$t_1$ is reachable from every vertex (for instance
$v_0\to v_3\to v_1\to t_1$) and there is no opponent, $w^\ast\equiv1$.

Write a Max strategy as the bit vector of its choices at
$(v_0,v_1,v_2,v_3)$, bit $0$ meaning the first listed successor. A
strategy is optimal exactly when its unique outgoing path from every
vertex reaches $t_1$. Reaching $t_1$ requires passing through $v_1$ with
$v_1\to t_1$, forcing bit $1$ to be $0$; and $v_1$ is reachable only via
$v_3\to v_1$, forcing bit $3$ to be $1$. Given those, $v_0\to v_3$ (bit
$0$ equal to $0$) succeeds for either choice at $v_2$, whereas
$v_0\to v_2$ succeeds only if $v_2\to v_3$. Hence
\[
  \Opt=\bigl\{(0,0,0,1),\;(0,0,1,1),\;(1,0,0,1)\bigr\},
\]
which exact rational computation of all $16$ value vectors confirms; the
one excluded completion, $(1,0,1,1)$, has
$\val_\sigma=(0,1,0,1)$. The coordinate projections of $\Opt$ are
$\{0,1\},\{0\},\{0,1\},\{1\}$, whose product has four elements while
$|\Opt|=3$. So $\Opt$ is not a subcube, and it is not even a sublattice:
$(0,0,1,1)\vee(1,0,0,1)=(1,0,1,1)\notin\Opt$.
\end{example}

\begin{example}[Greedy policy iteration for Min is unsound without stopping]\label{ex:pi-fails}
Two vertices suffice. Let $\Vmin=\{m\}$ and $\Vavg=\{c\}$ with
\[
  m\to\{t_1,\,c\},\qquad c\to\{m,\,m\},
\]
and $\Vmax=\emptyset$ (so $\sigma$ is the empty strategy). The true value
is $\val_\sigma\equiv0$: Min plays $m\to c$ and the token cycles
$m\to c\to m$ forever, never reaching $t_1$. Greedy policy iteration for
Min, started at $\tau(m)=t_1$, computes the values $\val(m)=\val(c)=1$
under that $\tau$, then evaluates the alternative successor $c$
\emph{against those same values}, finds $\val(c)=1\not<1=\val(t_1)$, and
therefore performs no switch --- returning $1$ instead of $0$.

The error is exactly the failure of \Cref{cor:comparison} in the absence
of stopping (\Cref{rem:no-postfixed}). Here $\Vmax=\emptyset$, so
$T_\sigma=T$; the vector $x\equiv1$ satisfies $Tx=x$, hence certainly
$x\le Tx$, yet $\val\equiv0$, so $x\le\val$ fails. In a stopping game
\Cref{cor:comparison} would force $x\le\val$ and no such fixed point
could exist. This is why every evaluation step below is performed by
linear programming or by explicit least-fixed-point computation rather
than by greedy improvement, and it is a further reason to pass to a
stopping game at the outset.
\end{example}

We now come to the main positive structural statement. It says that
improving paths are always \emph{short}; the difficulty of
\Cref{prob:main} is therefore entirely one of \emph{finding} them.

\begin{lemma}[Single-switch improvement]\label{lem:switch}
Let $G$ be stopping, $\sigma$ a Max strategy and $v\in S_\sigma$. Then
$\val_{\sigma[v]}\ge\val_\sigma$ pointwise, with strict inequality at $v$.
Consequently $\Phi(\val_{\sigma[v]})>\Phi(\val_\sigma)$, where
$\Phi(x):=\sum_{u\in V}x(u)$.
\end{lemma}

\begin{proof}
Put $x:=\val_\sigma$ and $\sigma':=\sigma[v]$. Then
$(T_{\sigma'}x)(v)=x(\bar\sigma(v))>x(\sigma(v))=x(v)$, while
$(T_{\sigma'}x)(u)=(T_\sigma x)(u)=x(u)$ for every other $u$; hence
$x\le T_{\sigma'}x$ with strict inequality at $v$. By
\Cref{cor:comparison} --- this is where stopping is used ---
$x\le\val_{\sigma'}$. Applying $T_{\sigma'}$ once more at the coordinate
$v$ and using monotonicity,
$\val_{\sigma'}(v)=(T_{\sigma'}\val_{\sigma'})(v)\ge(T_{\sigma'}x)(v)>x(v)$.
The statement about $\Phi$ follows because $\Phi$ is strictly increasing
on pointwise-comparable distinct vectors.
\end{proof}

\begin{theorem}[Short improving paths]\label{thm:short-path}
Let $G$ be a stopping SSG, let $\sigma\notin\Opt$ and let
$\sigma^\ast\in\Opt$ be arbitrary. Then there exists a vertex
$v\in\Vmax$ with $\sigma(v)\ne\sigma^\ast(v)$ that is strictly
switchable at $\sigma$. Consequently, writing
$d(\sigma):=|\{v:\sigma(v)\ne\sigma^\ast(v)\}|$, there is a sequence
\[
  \sigma=\sigma_0,\ \sigma_1,\dots,\ \sigma_r\in\Opt,
  \qquad r\le d(\sigma)\le n=|\Vmax|,
\]
in which each $\sigma_{j+1}=\sigma_j[v_j]$ is obtained by a single
strictly switchable flip and $\val_{\sigma_0}\lneq\val_{\sigma_1}\lneq\cdots\lneq\val_{\sigma_r}=w^\ast$.
\end{theorem}

\begin{proof}
Suppose, for contradiction, that no vertex of
$D:=\{v:\sigma(v)\ne\sigma^\ast(v)\}$ is strictly switchable at $\sigma$;
that is, $\val_\sigma(\sigma^\ast(v))\le\val_\sigma(\sigma(v))$ for every
$v\in D$ (note $\bar\sigma(v)=\sigma^\ast(v)$ for $v\in D$, because each
vertex has exactly two successors). Put $x:=\val_\sigma$. For
$v\in\Vmax\setminus D$ we have $\sigma^\ast(v)=\sigma(v)$ and hence
$(T_{\sigma^\ast}x)(v)=x(\sigma(v))=x(v)$; for $v\in D$,
$(T_{\sigma^\ast}x)(v)=x(\sigma^\ast(v))\le x(\sigma(v))=x(v)$. At
$\Vmin$ and $\Vavg$ vertices and at the sinks the operators
$T_{\sigma^\ast}$ and $T_\sigma$ agree, and $x=T_\sigma x$, so equality
holds there. Therefore $T_{\sigma^\ast}x\le x$, and \Cref{cor:comparison}
gives $\val_{\sigma^\ast}\le x=\val_\sigma$. Since
$\val_{\sigma^\ast}=w^\ast$ and $\val_\sigma\le w^\ast$ always, we get
$\val_\sigma=w^\ast$, i.e.\ $\sigma\in\Opt$ --- contradicting the
hypothesis. This proves the first assertion.

For the second, note that flipping $v\in D$ turns $\sigma(v)$ into
$\sigma^\ast(v)$, so $d$ strictly decreases at each step, and by
\Cref{lem:switch} each such flip strictly improves the value vector.
Iterating at most
$d(\sigma)\le n$ times, and stopping as soon as the current strategy lies
in $\Opt$, produces the asserted chain; it terminates in $\Opt$ because
if $\sigma_j\notin\Opt$ the first assertion supplies a further flip
inside the (shrinking) disagreement set.
\end{proof}

\Cref{thm:short-path} says a wrong switchable vertex exists somewhere. It
can in fact be localised, to the vertices at which $\sigma$ loses the most
value.

\begin{lemma}[Maximum-deficit localisation]\label{lem:max-deficit}
Let $G$ be stopping and let $\sigma\notin\Opt$. Put $x:=\val_\sigma$,
$g:=w^\ast-x\ge0$, $M:=\max_u g(u)>0$ and $Z:=\{u\in V:g(u)=M\}$. Then
$Z$ contains a Max vertex $v$ that is both wrong at $\sigma$ and strictly
switchable at $\sigma$.
\end{lemma}

\begin{proof}
First, $Z$ contains no sink, since $g$ vanishes there and $M>0$. We
examine what a vertex of $Z$ can do.

If $u\in Z\cap\Vavg$ then $g(u)=\tfrac12(g(u^{(0)})+g(u^{(1)}))$, and
since both terms are at most $M$, \emph{both} successors lie in $Z$.

If $u\in Z\cap\Vmin$, pick $i$ with $x(u)=x(u^{(i)})$. Then
$w^\ast(u)\le w^\ast(u^{(i)})$, so
$M=g(u)=w^\ast(u)-x(u^{(i)})\le w^\ast(u^{(i)})-x(u^{(i)})=g(u^{(i)})\le M$;
hence $u^{(i)}\in Z$.

If $u\in Z\cap\Vmax$ is \emph{not} wrong, then $w^\ast(\sigma(u))=w^\ast(u)$
and $x(\sigma(u))=x(u)$, so $g(\sigma(u))=g(u)=M$ and $\sigma(u)\in Z$.

If $u\in Z\cap\Vmax$ \emph{is} wrong and moreover
$x(\bar\sigma(u))=x(\sigma(u))=x(u)$, then $w^\ast(\bar\sigma(u))=w^\ast(u)$
(the other successor attains the maximum), so
$g(\bar\sigma(u))=w^\ast(u)-x(u)=M$ and $\bar\sigma(u)\in Z$.

Suppose now, for contradiction, that no vertex of $Z\cap\Vmax$ is both
wrong and strictly switchable. Note first that a wrong $v\in Z$ always
satisfies $x(\bar\sigma(v))\ge x(\sigma(v))$: indeed
$g(\bar\sigma(v))\le M=g(v)=w^\ast(\bar\sigma(v))-x(v)$ gives
$x(\bar\sigma(v))\ge x(v)=x(\sigma(v))$. By assumption the inequality is
never strict on $Z$, so every wrong $v\in Z$ falls into the fourth case
above. Define a Max strategy $\sigma'$ agreeing with $\sigma$ except that
$\sigma'(v):=\bar\sigma(v)$ at every wrong $v\in Z$, and a Min strategy
$\tau'$ choosing at each $u\in Z\cap\Vmin$ a successor in $Z$ (and
arbitrarily elsewhere). By the four cases, every vertex of $Z$ then moves
inside $Z$ with probability $1$. So under $(\sigma',\tau')$ a token
started in $Z$ never leaves $Z$, and since $Z$ contains no sink it is
never absorbed --- contradicting the assumption that $G$ is stopping.
\end{proof}

\begin{remark}
\Cref{lem:max-deficit} is sharp in a way worth recording: a wrong vertex
of $Z$ need \emph{not} be strictly switchable --- among $2976$
(game, non-optimal strategy) pairs drawn from stopping games, $481$ wrong
vertices of $Z$ had $x(\bar\sigma(v))=x(\sigma(v))$ --- but in every one
of those $2976$ pairs at least one wrong vertex of $Z$ was strictly
switchable, as the lemma asserts. The lemma does not yield an algorithm:
$Z$ is defined through $w^\ast$, so computing it is again the whole
problem. What it does is identify \emph{which} vertices a selection rule
should be trying to find.
\end{remark}

\begin{corollary}[The difficulty is selection, not distance]\label{cor:selection}
For every stopping SSG and every Max strategy $\sigma$, the improving
switch digraph contains a directed path of length at most $n$ from
$\sigma$ to an optimal strategy. Hence no superpolynomial lower bound on
any strategy-improvement algorithm can arise from the \emph{length} of
improving paths: such a path of length at most $n$ always exists, so any
superpolynomial behaviour is a property of the algorithm's
\emph{selection rule} alone. (This says nothing about lower bounds for
\Cref{prob:main} by other means.)
\end{corollary}

\begin{proof}
Immediate from \Cref{thm:short-path}.
\end{proof}

\Cref{thm:short-path} guarantees that a short improving path exists; it
does not say that a natural rule follows one. The most natural rule does
not, even in a stopping game.

\begin{proposition}[All-switches overshoots]\label{prop:allswitch-overshoot}
There is a stopping SSG with $|\Vmax|=14$ on which the all-switches rule
--- from a strategy $\sigma$, flip \emph{every} vertex of $S_\sigma$ ---
performs $16>|\Vmax|$ iterations before reaching an optimal strategy,
switching four Max vertices three times each and four more twice each.
Explicitly, let $B$ be the game on non-sink vertices $0,\dots,27$ with
$t_0=28$, $t_1=29$, types
\[
\begin{aligned}
 &(\max,\max,\mathrm{avg},\mathrm{avg},\mathrm{avg},\mathrm{avg},\mathrm{avg},\mathrm{avg},\mathrm{avg},\max,\mathrm{avg},\mathrm{avg},\mathrm{avg},\max,\\
 &\ \ \max,\mathrm{avg},\mathrm{avg},\max,\max,\mathrm{avg},\max,\mathrm{avg},\max,\max,\max,\max,\max,\max)
\end{aligned}
\]
and out-edge pairs, in vertex order,
\[
\begin{aligned}
 &(9,26),(11,20),(11,28),(21,10),(13,22),(19,23),(5,28),(27,29),(25,20),(25,10),\\
 &(2,17),(14,26),(8,6),(26,15),(27,9),(29,9),(0,10),(25,3),(6,13),(12,25),\\
 &(16,19),(1,6),(27,24),(24,18),(14,1),(25,14),(10,0),(14,22),
\end{aligned}
\]
and let $G:=B_7$ be its damping (\Cref{def:damping}) with $m=7$, a
stopping game on $422$ vertices ($420$ of them non-sink) with the same
$14$ Max vertices. Started
from the strategy taking every first-listed edge, all-switches performs
$16$ iterations; vertices $1,20,22,24$ are switched three times and
$0,9,14,17$ twice; the value vectors increase strictly at every step.
\end{proposition}

\begin{proof}
$G$ is stopping by \Cref{lem:gadget}. The run is a finite deterministic
computation: each iteration evaluates $\val_{\sigma}$ by solving one
linear system with rational coefficients (there are no Min vertices, so
$\val_\sigma=v_{\sigma,\cdot}$ is immediate from \Cref{lem:lfp}), and
then flips every strictly switchable vertex. Carrying it out in exact
rational arithmetic yields the sizes of the successive switched sets,
namely $1,1,1,2,4,2,5,1,1,1,1,1,1,2,1,1$, the switch multiset
$\{3,3,3,3,2,2,2,2,1,1,1,1,1,1\}$, and a strictly increasing sequence of
$17$ distinct value vectors, terminating with $S_\sigma=\emptyset$; the
final value at vertex $0$ is a rational with a $36$-digit numerator and a
$37$-digit denominator, approximately $0.133044$.
\end{proof}

\begin{remark}
The contrast with \Cref{thm:short-path} is the point. A path of at most
$14$ single improving switches to $\Opt$ exists from the same starting
strategy; all-switches instead takes $16$ steps and \emph{undoes} its own
earlier choices, flipping four vertices three times. Combined with
\Cref{thm:all-switches-refuted}, which shows the all-switches successor
can be strictly dominated by a single-switch successor, this rules out
all-switches as a candidate for \textsc{(Poly-Rule)} on any argument that
counts vertices or measures progress towards the optimum vertex by
vertex.
\end{remark}

\subsection{Exactly how much a single switch gains}

The switching lemma inside \Cref{thm:short-path} is qualitative: a
strictly switchable flip improves the value. The improvement can be
written down exactly, and the formula shows precisely why no potential of
polynomial range can control strategy iteration.

For a fixed positional pair $(\sigma,\tau)$ of a stopping game write $P$
for the transition matrix of $M_{\sigma,\tau}$, $f:=v_{\sigma,\tau}$, and
for vertices $x,u$ let
\[
  h(x\to u)\;:=\;\Pr\nolimits^{P}_{x}\bigl[\text{the token ever visits } u\bigr],
  \qquad h(u\to u):=1 .
\]

\begin{theorem}[Transfer-impedance formula]\label{thm:impedance}
Let $G$ be stopping, $(\sigma,\tau)$ a positional pair with chain $P$ and
value $f$, and let $u$ be a controlled vertex (in $\Vmax$ or $\Vmin$)
whose chosen successor is $a$; let $b$ be its other successor and let
$P'$ be the chain obtained by switching $u$ from $a$ to $b$, with value
$f'$. Then $h(\cdot\to u)$ is the same function for $P$ and $P'$,
$h(b\to u)<1$, and for every vertex $v$
\[
  f'(v)-f(v)\;=\;\bigl(f(b)-f(a)\bigr)\,\frac{h(v\to u)}{1-h(b\to u)} .
\]
\end{theorem}

\begin{proof}
\emph{$h$ is unchanged.} For $x\ne u$, first-step analysis gives
$h(x\to u)=\sum_w P(x,w)h(w\to u)$ with $h(u\to u)=1$ and
$h(\text{sink}\to u)=0$; this linear system involves only the rows of $P$
indexed by $x\ne u$, and $P$ and $P'$ agree on exactly those rows. The
system has a unique solution because the substochastic matrix obtained
from $P$ by deleting row and column $u$ has spectral radius below $1$
(from every vertex a sink is reached almost surely, by
\Cref{lem:absorption}). Hence $h$ is the same for $P$ and $P'$.

\emph{$h(b\to u)<1$.} In $P'$ the token moves from $u$ to $b$ with
probability $1$, so the probability of returning to $u$ from $u$ equals
$h(b\to u)$. If this were $1$, then $u$ would be revisited infinitely
often almost surely and the token would never be absorbed, contradicting
that $G$ is stopping.

\emph{The formula.} Both $f$ and $f'$ solve their own systems:
$f=Pf+\beta$ and $f'=P'f'+\beta'$, where $\beta$ (resp.\ $\beta'$) is the
vector of one-step probabilities of landing in $t_1$. Subtracting,
\[
  f'-f=P'(f'-f)+\Delta,\qquad \Delta:=(P'f+\beta')-f ,
\]
so $(I-P')(f'-f)=\Delta$ and $f'-f=\mathcal{N}'\Delta$ with
$\mathcal{N}':=(I-P')^{-1}=\sum_{k\ge0}(P')^{k}$, the Green function,
which is finite because $P'$ has spectral radius below $1$. Now $P'$ and
$P$ agree off row $u$, and so do $\beta',\beta$; hence $\Delta$ is
supported on $\{u\}$, where
$\Delta(u)=f(b)-f(u)=f(b)-f(a)$ using $f(u)=f(a)$. Therefore
$f'-f=\bigl(f(b)-f(a)\bigr)\,\mathcal{N}'(\cdot,u)$. Finally
$\mathcal{N}'(v,u)=\mathbb{E}_v[\#\text{visits to }u]$ decomposes at the
first visit as $h'(v\to u)\,\mathcal{N}'(u,u)$, and
$\mathcal{N}'(u,u)=\bigl(1-h(b\to u)\bigr)^{-1}$ by the previous
paragraph; since $h'=h$, the claim follows.
\end{proof}

\begin{remark}\label{rem:impedance}
The identity was verified in exact rational arithmetic on $544$
(game, profile, switched vertex) instances over stopping games, with no
mismatch, and the equality $h=h'$ was confirmed in every case.

Its content is this. The numerator $h(v\to u)\in[0,1]$ is a bounded
\emph{transfer impedance}: it says how strongly a perturbation at $u$ is
felt at $v$, and it vanishes for vertices that cannot reach $u$. The
denominator $1-h(b\to u)$ is an \emph{escape probability}, and nothing
bounds it away from $0$. It can be as small as $2^{-\Theta(N)}$: take
$u\in\Vmax$ with $u\to\{t_1,c_1\}$ and a chain of $k$ average vertices
$c_1,\dots,c_k$ with $c_i\to\{u,c_{i+1}\}$ for $i<k$ and
$c_k\to\{u,t_0\}$. Escaping $u$ from $c_1$ requires descending the whole
chain, so $h(c_1\to u)=1-2^{-k}$ and the escape probability is exactly
$2^{-k}$ on $k+1$ non-sink vertices (verified in exact arithmetic for
$k=2,4,6,8,10$). Consequently the gain of a switch is not a function of
local data of polynomial range, and any argument that bounds the number of
iterations by
$(\text{range of a potential})/(\text{minimum increment})$ is
inherently exponential --- independently of which potential is chosen.
\end{remark}

\section{Four refutations}\label{sec:refutations}

The results of this section eliminate four specific and otherwise
natural routes to a polynomial bound. Each is witnessed by an explicit
game whose values were computed in exact rational arithmetic.

\subsection{The all-switches rule does not dominate}

A natural engine for transferring counting arguments from
one-player Markov decision processes to SSGs is the assertion that
switching \emph{all} strictly switchable vertices makes at least as much
progress as switching any subset. It is false.

\begin{theorem}[All-switches dominance fails]\label{thm:all-switches-refuted}
There is a stopping SSG $G$, with five non-sink vertices, and a Max
strategy $\sigma$ with $|S_\sigma|=2$, such that for the single-vertex
switch $U=\{x\}\subsetneq S_\sigma$,
\[
  \val_{\sigma[U]}\;\gneq\;\val_{\sigma[S_\sigma]} \qquad\text{pointwise.}
\]
Explicitly, let $\Vmax=\{x,y\}$, $\Vmin=\{m\}$, $\Vavg=\{a,h\}$ and
\[
  x\to\{t_0,a\},\quad y\to\{m,h\},\quad m\to\{a,x\},\quad a\to\{y,t_1\},\quad h\to\{t_1,t_0\},
\]
with $\sigma=(x\mapsto t_0,\ y\mapsto m)$. Then
$S_\sigma=\{x,y\}$ and, ordering coordinates $(x,y,m,a,h)$,
\[
  \val_\sigma=(0,0,0,\tfrac12,\tfrac12),\qquad
  \val_{\sigma[\{x\}]}=(1,1,1,1,\tfrac12),\qquad
  \val_{\sigma[\{x,y\}]}=(\tfrac34,\tfrac12,\tfrac34,\tfrac34,\tfrac12).
\]
Moreover $\val_{\sigma[\{x\}]}$ is the value vector $w^\ast$ of $G$, so a
single switch reaches the optimum while the all-switches successor does
not.
\end{theorem}

\begin{proof}
The game is stopping: the average vertices $a$ and $h$ each have an edge
to a sink, and from every vertex, under every strategy pair, one of $a$,
$h$ or a sink is reachable, so \Cref{def:stopping} holds. The three value
vectors are obtained by solving the three induced systems. For
$\sigma=(x\mapsto t_0, y\mapsto m)$: Min's best response is
$m\mapsto x$, and then $x\to t_0$ gives $\val(x)=0$, hence
$\val(m)=0$, $\val(y)=\val(m)=0$; and $a\to\{y,t_1\}$ gives
$\val(a)=\tfrac12(0+1)=\tfrac12$, $\val(h)=\tfrac12$. Both $x$ and $y$
are strictly switchable: $\val_\sigma(a)=\tfrac12>0=\val_\sigma(t_0)$ and
$\val_\sigma(h)=\tfrac12>0=\val_\sigma(m)$.

For $\sigma[\{x\}]=(x\mapsto a,\ y\mapsto m)$: whichever successor Min
selects at $m$, the token from $a$ reaches $t_1$ with probability
$\tfrac12$ immediately and otherwise moves to $y\to m\to\{a,x\}$, and
$x\to a$; thus the token returns to $a$ with probability $1$ whenever it
does not stop at $t_1$, so $t_1$ is reached with probability $1$ from
each of $x,y,m,a$, giving $(1,1,1,1,\tfrac12)$. As this vector satisfies
$Tw=w$ and $G$ is stopping, it equals $w^\ast$ by
\Cref{thm:contraction}.

For $\sigma[\{x,y\}]=(x\mapsto a,\ y\mapsto h)$: Min's best response at
$m$ is $m\mapsto x$, and the system
$\val(a)=\tfrac12(\val(y)+1)$, $\val(y)=\val(h)=\tfrac12$,
$\val(x)=\val(a)$, $\val(m)=\val(x)$ yields
$\val(a)=\tfrac12(\tfrac12+1)=\tfrac34$ and hence
$(\tfrac34,\tfrac12,\tfrac34,\tfrac34,\tfrac12)$. Comparing coordinatewise
with $(1,1,1,1,\tfrac12)$ establishes the strict domination.
\end{proof}

\begin{corollary}
No proof of a polynomial iteration bound for the all-switches rule can
proceed by showing that its successor dominates the successors of all
partial switches; and the all-switches rule may require strictly more
iterations from a given position than a best single switch.
\end{corollary}

\subsection{The induced orientation of the two-player cube can be cyclic}

On the Max-only cube the map $\sigma\mapsto\val_\sigma$ is a strictly
increasing potential along improving switches (\Cref{lem:switch}), so
that orientation is acyclic. It is tempting to work instead on the
\emph{combined} cube of pairs $(\sigma,\tau)$, orienting each edge in the
direction its owner prefers; algorithms and lower-bound machinery for
unique-sink orientations are frequently stated under an additional
\emph{acyclicity} hypothesis. The following shows that hypothesis is
unavailable on the combined cube.

We orient the cube $\{0,1\}^{\Vmax\cup\Vmin}$ of positional profiles by
directing the edge in coordinate $v$ away from a profile exactly when
switching $v$ is improving for its owner --- increasing $v_{\sigma,\tau}(v)$
if $v\in\Vmax$, decreasing it if $v\in\Vmin$. We make no claim here about
any linear-complementarity formulation; the statement below concerns this
improving-switch orientation only.

\begin{theorem}[The combined-cube orientation can be cyclic]\label{thm:cyclic-uso}
There is a stopping game with two Max vertices $a,b$ and one Min vertex
$m$ whose induced orientation of the $3$-cube of strategy profiles is a
unique-sink orientation containing a directed $6$-cycle. Explicitly,
taking each option to be a sub-probability distribution over
$\{a,b,m\}$ together with a weight on $t_1$ (all dyadic with denominator
$16$, hence realisable by depth-$4$ trees of average vertices):
\[
\begin{array}{llll}
  a: &\text{option }0:\ \tfrac12\!\to\!a,\ \tfrac18\!\to\!t_1;
     &\text{option }1:\ \tfrac78\!\to\!m,\ \tfrac1{16}\!\to\!t_1;\\
  b: &\text{option }0:\ \tfrac{15}{16}\!\to\!a;
     &\text{option }1:\ \tfrac12\!\to\!b,\ \tfrac1{16}\!\to\!t_1;\\
  m: &\text{option }0:\ \tfrac12\!\to\!a;
     &\text{option }1:\ \tfrac{15}{16}\!\to\!b.
\end{array}
\]
Every option has total mass at most $\tfrac{15}{16}<1$, so the game is
stopping. Orienting, at each profile $S\in\{0,1\}^{\{a,b,m\}}$, the edge
in coordinate $v$ away from $S$ exactly when switching $v$ is improving
for its owner, one obtains a USO --- all $27$ faces have a unique sink
and no comparison is tied --- whose vertex values are
\[
\begin{array}{ll}
  S=(0,0,0):\ \bigl(\tfrac14,\tfrac{15}{64},\tfrac18\bigr),\ \text{out}=\emptyset;
  & S=(0,0,1):\ \bigl(\tfrac14,\tfrac{15}{64},\tfrac{225}{1024}\bigr),\ \text{out}=\{a,m\};\\
  S=(0,1,0):\ \bigl(\tfrac14,\tfrac18,\tfrac18\bigr),\ \text{out}=\{b,m\};
  & S=(0,1,1):\ \bigl(\tfrac14,\tfrac18,\tfrac{15}{128}\bigr),\ \text{out}=\{b\};\\
  S=(1,0,0):\ \bigl(\tfrac19,\tfrac5{48},\tfrac1{18}\bigr),\ \text{out}=\{a,b\};
  & S=(1,0,1):\ \bigl(\tfrac{128}{473},\tfrac{120}{473},\tfrac{225}{946}\bigr),\ \text{out}=\{m\};\\
  S=(1,1,0):\ \bigl(\tfrac19,\tfrac18,\tfrac1{18}\bigr),\ \text{out}=\{a\};
  & S=(1,1,1):\ \bigl(\tfrac{169}{1024},\tfrac18,\tfrac{15}{128}\bigr),\ \text{out}=\{a,b,m\}.
\end{array}
\]
The six profiles other than the sink $(0,0,0)$ and the source $(1,1,1)$
form the directed cycle
\[
  (0,0,1)\to(1,0,1)\to(1,0,0)\to(1,1,0)\to(0,1,0)\to(0,1,1)\to(0,0,1),
\]
with flip sequence $a,m,b,a,m,b$.
\end{theorem}

\begin{proof}
Each of the eight profiles induces a Markov chain whose value vector is
the unique solution of a $3\times3$ linear system with rational
coefficients; solving the eight systems exactly gives the table. Each of
the $24$ vertex--coordinate comparisons is then a comparison of two
explicit rationals; all $24$ are strict, so the orientation is
well defined and degeneracy does not arise. For instance at $S=(0,0,1)$
and coordinate $a$: the current value is $\tfrac14$, while option $1$
evaluates to $\tfrac1{16}+\tfrac78\cdot\tfrac{225}{1024}=\tfrac{2087}{8192}>\tfrac14$,
so the $a$-edge points away from $S$. Checking each of the $27$ faces of
the $3$-cube (the cube itself, six facets, twelve edges, eight vertices)
exhibits exactly one sink in each, so the orientation is a USO. Finally,
each of the six listed transitions flips a coordinate lying in the
out-set of its tail, so the displayed closed walk is a directed cycle.
\end{proof}

\begin{corollary}\label{cor:no-potential}
There is no real-valued potential on the profiles of a general SSG that
strictly increases along every improving edge of the two-player strategy
cube. In particular, algorithmic and lower-bound machinery that assumes
acyclic unique-sink orientations does not apply to the combined
two-player cube.
\end{corollary}

\begin{proof}
A directed cycle forbids a strictly increasing potential; the cycle is
supplied by \Cref{thm:cyclic-uso}.
\end{proof}

\begin{remark}
\Cref{cor:no-potential} does not contradict \Cref{thm:short-path}: the
latter concerns the \emph{Max-only} cube, on which $\sigma\mapsto\val_\sigma$
is a strictly increasing potential along improving switches by
\Cref{cor:comparison}. Acyclicity holds on each one-player cube and fails
on the combined cube.
\end{remark}

\subsection{Value iteration is exponential already without players}

Iterating $T$ from $0$ converges to $\val$ (\Cref{thm:contraction}), and
by \Cref{lem:denominator} it suffices to approximate $\val(v_0)$ to
within $1/(2D)$ to decide \Cref{prob:main}. The following family shows
that the number of iterations required is exponential, and that this has
nothing to do with the presence of players: the family has none.

\begin{theorem}[Exponential lower bound for value iteration]\label{thm:vi-lower}
For $m\ge3$ let $G_m$ be the SSG with non-sink vertices
$c_1,\dots,c_L$, $L=2m$, \emph{all} in $\Vavg$ (so
$\Vmax=\Vmin=\emptyset$), and edges
\[
  c_i\to\{o_i,\;c_{i+1}\}\ \ (1\le i\le L-1),\qquad c_L\to\{t_0,\;t_1\},
\]
where $o_i=c_1$ for $1\le i\le m-1$, $o_m=t_1$, and $o_i=t_0$ for
$m+1\le i\le 2m-1$. Then $G_m$ is stopping, has $2m$ non-sink vertices,
and
\begin{enumerate}[label=(\roman*),leftmargin=2.2em]
\item $\val(c_1)=\tfrac12+2^{-(m+1)}$, so $\val(c_1)>\tfrac12$ by an
  exponentially small margin;
\item $\min_i\val(c_i)=2^{-m}$;
\item $\bigl(T^{k}0\bigr)(c_1)<\tfrac12$ for every
  $k<\ln 2\cdot 2^{m-2}$.
\end{enumerate}
Consequently value iteration from $0$ requires $2^{\Omega(N)}$ iterations
before its iterate at $c_1$ can certify $\val(c_1)\ge\tfrac12$.
\end{theorem}

\begin{proof}
$G_m$ is stopping because from every $c_i$ the path
$c_i\to c_{i+1}\to\cdots\to c_L\to t_0$ has positive probability.
As there are no players, $T$ is the affine map $x\mapsto Mx+b$ with $M$
the substochastic transition matrix and $b$ the vector of one-step
probabilities of landing in $t_1$, and $\val$ is its unique fixed point.

(i) Call the segment of the trajectory from $c_1$ until the token is
first at a sink or back at $c_1$ an \emph{excursion}. The token reaches
$c_i$ without leaving earlier with probability $2^{-(i-1)}$ and then
exits to $o_i$ with probability $\tfrac12$, so it exits at $o_i$ with
probability $2^{-i}$ for $i\le L-1$, and reaches each of $t_0,t_1$ from
$c_L$ with probability $2^{-L}$. Hence the excursion ends at $t_1$ with
probability $p=2^{-m}+2^{-2m}$, at $t_0$ with probability
$q=\sum_{i=m+1}^{2m-1}2^{-i}+2^{-2m}=2^{-m}-2^{-2m}$, and returns to
$c_1$ with probability $\sum_{i=1}^{m-1}2^{-i}=1-2^{-(m-1)}$; these sum
to $1$. By the strong Markov property
$\val(c_1)=p+(1-p-q)\val(c_1)$, so
$\val(c_1)=p/(p+q)=(2^{-m}+2^{-2m})/2^{-m+1}=\tfrac12+2^{-(m+1)}$.

(ii) For $i\ge m+1$, $\val(c_i)=\tfrac12\val(c_{i+1})$ with
$\val(c_L)=\tfrac12$, giving $\val(c_i)=2^{-(2m-i+1)}$, least at $i=m+1$
where it equals $2^{-m}$. For $i\le m-1$,
$\val(c_i)=\tfrac12\val(c_1)+\tfrac12\val(c_{i+1})\ge\tfrac14$, and
$\val(c_m)=\tfrac12+\tfrac12\val(c_{m+1})$. Hence the minimum is
$2^{-m}$.

(iii) Let $e^{(k)}:=\val-T^{k}0$. Since $T$ is affine with fixed point
$\val$, $e^{(k)}=M^{k}\val$, and as $M\ge0$ and $\val\ge0$,
\[
  e^{(k)}(c_1)\;=\;\mathbb{E}_{c_1}\bigl[\val(X_k)\mathbf{1}\{X_k\ \text{not absorbed}\}\bigr]
  \;\ge\;\bigl(\min_i\val(c_i)\bigr)\Pr\nolimits_{c_1}[\text{alive at }k]
  \;\ge\;2^{-m}\bigl(1-2^{1-m}\bigr)^{k},
\]
the last step because each excursion is absorbed with probability
$p+q=2^{1-m}$, independently across excursions, and each excursion takes
at least one step, so surviving $k$ excursions implies surviving $k$
steps. Now $\val(c_1)-\tfrac12=2^{-(m+1)}$, so
$\bigl(T^{k}0\bigr)(c_1)<\tfrac12$ whenever
$e^{(k)}(c_1)>2^{-(m+1)}$, i.e.\ whenever $(1-2^{1-m})^{k}>\tfrac12$.
Since $-\ln(1-u)\le 2u$ for $0<u\le\tfrac12$, taking $u=2^{1-m}$
(legitimate for $m\ge3$) this holds for every
$k<\ln2/(2\cdot2^{1-m})=\ln2\cdot2^{m-2}$.
\end{proof}

\begin{remark}
Exact rational computation confirms
$\val(c_1)=\tfrac12+2^{-(m+1)}$ and $\min_i\val(c_i)=2^{-m}$ for
$m=3,\dots,12$, and the number of iterations before
$\bigl(T^{k}0\bigr)(c_1)$ first reaches $\tfrac12$ is measured to be
$13,37,99,248,595,1385,3149,7041,15545,33985$ for those $m$ --- growing
by a factor close to $2.2$ per increment of $m$, comfortably above the
proved bound $\ln2\cdot 2^{m-2}$.

Two consequences deserve emphasis. First, the obstruction is not the
\emph{precision} of the arithmetic but the \emph{mixing time} of the
chain: the margin $\val(c_1)-\tfrac12=2^{-(m+1)}$ is exponentially small,
so no fixed polynomial number of iterations at any precision can decide
the instance. Second, since $G_m$ has no Max or Min vertices at all, the
failure of value iteration is not caused by the interaction of the two
players; it is already present in the one-dimensional problem of
computing an absorption probability. Any polynomial algorithm must
therefore compute values by solving the linear system exactly rather than
by iterating $T$.
\end{remark}

\subsection{Improving switches need not point toward the optimal set}

By \Cref{thm:opt-subcube}, in a stopping game $\Opt=\prod_{v}A_v$ where
$A_v\subseteq\{0,1\}$ is the set of choices at $v$ attaining
$\max\bigl(w^\ast(v^{(0)}),w^\ast(v^{(1)})\bigr)$. This makes the
Hamming distance to the optimal set explicit:
\[
  d(\sigma)\;:=\;\bigl|\{v\in\Vmax : \sigma(v)\notin A_v\}\bigr|
  \;=\;\min_{\sigma^\ast\in\Opt}\bigl|\{v:\sigma(v)\ne\sigma^\ast(v)\}\bigr| .
\]
If every strictly switchable vertex were a vertex at which $\sigma$ is
\emph{wrong}, i.e.\ if $S_\sigma\subseteq\{v:\sigma(v)\notin A_v\}$ always
held, then \emph{every} single improving switch would decrease $d$ by
one, and strategy iteration under \emph{any} selection rule would
terminate within $n$ steps --- giving a polynomial algorithm. It does
not hold.

\begin{theorem}[No naive Hamming potential]\label{thm:hamming-refuted}
There is a stopping SSG with two Max vertices, no Min vertices and no
average vertices, and a Max strategy $\sigma$, such that some
$v\in S_\sigma$ satisfies $\sigma(v)\in A_v$; switching $v$ is strictly
improving yet leaves $d(\sigma)$ unchanged. Explicitly, let
$\Vmax=\{p,q\}$ with
\[
  p\to\{t_1,\,t_0\},\qquad q\to\{t_1,\,p\},
\]
and let $\sigma$ be $p\mapsto t_0$, $q\mapsto p$. Then $w^\ast=(1,1)$,
$A_p=\{t_1\}$, $A_q=\{t_1,p\}$, $\Opt=\{p\mapsto t_1\}\times\{t_1,p\}$,
and $\val_\sigma=(0,0)$. Both $p$ and $q$ are strictly switchable at
$\sigma$, but $\sigma(q)=p\in A_q$, so switching $q$ yields
$\val=(0,1)$ --- a strict improvement --- while $d$ remains $1$.
\end{theorem}

\begin{proof}
The game is stopping: $p$ moves to a sink in one step under either
choice, and $q$ moves to a sink or to $p$. Its value is $w^\ast=(1,1)$,
since Max may play $p\to t_1$ and $q\to t_1$. Hence
$w^\ast(p)=w^\ast(t_1)=1$, so both choices at $q$ attain the maximum and
$A_q=\{t_1,p\}$, while $w^\ast(t_1)=1>0=w^\ast(t_0)$ gives
$A_p=\{t_1\}$. Under $\sigma$ the token goes $q\to p\to t_0$, so
$\val_\sigma=(0,0)$; both switches are strictly improving because
$\val_\sigma(t_1)=1>0$. Finally $d(\sigma)=1$ (only $p$ is wrong) and
$d(\sigma[q])=1$ as well.
\end{proof}

\begin{corollary}\label{cor:hamming}
No proof of a polynomial iteration bound can proceed by arguing that an
arbitrary improving switch decreases the Hamming distance to $\Opt$. Any
Hamming-based argument must additionally explain \emph{which} switch to
take, which is the selection problem of \Cref{cor:selection}.
\end{corollary}

\begin{proof}
Immediate from \Cref{thm:hamming-refuted}: at the strategy exhibited
there the improving switch at $q$ does not reduce $d$, so $d$ is not
decreased by every improving switch.
\end{proof}

One might hope that although \emph{some} improving switch fails to be
wrong, a specific polynomial-time rule always selects a wrong one; by
\Cref{thm:short-path} a wrong switchable vertex always exists, so such a
rule would terminate in at most $n$ iterations and settle
\Cref{prob:main}. The following shows that the natural candidates do not
have this property.

\begin{proposition}\label{prop:rules-fail}
None of the following selection rules always chooses a wrong vertex from
$S_\sigma$: largest improvement gap
$\val_\sigma(\bar\sigma(v))-\val_\sigma(\sigma(v))$; smallest such gap;
largest $\val_\sigma(\bar\sigma(v))$; smallest $\val_\sigma(\sigma(v))$;
first or last index; and the one-step lookahead rule maximising
$\Phi(\val_{\sigma[v]})$, where $\Phi(x)=\sum_u x(u)$.
\end{proposition}

\begin{proof}
A single stopping game refutes the lookahead rule, the strongest of the
seven. Let $\Vmax=\{u_0,u_1,u_3,u_4\}$, $\Vavg=\{u_2\}$, with
\[
  u_0\to\{t_0,t_1\},\quad u_1\to\{t_0,u_0\},\quad u_2\to\{u_4,t_1\},\quad
  u_3\to\{t_0,u_1\},\quad u_4\to\{u_3,u_0\}.
\]
Then $w^\ast\equiv1$, so the greedy sets are $A_{u_0}=A_{u_1}=A_{u_3}=$
the second listed successor, while $A_{u_4}=\{u_3,u_0\}$ contains
\emph{both} choices (both have $w^\ast=1$); hence $u_4$ is never wrong.
Take $\sigma$ given by $u_0\mapsto t_1$, $u_1\mapsto t_0$,
$u_3\mapsto t_0$, $u_4\mapsto u_3$. Exact computation gives
$\val_\sigma=(1,0,\tfrac12,0,0)$ on $(u_0,\dots,u_4)$. Exactly two
vertices are strictly switchable: $u_1$, which is wrong, with
$\Phi(\val_{\sigma[u_1]})=\tfrac72$; and $u_4$, which is not wrong, with
$\Phi(\val_{\sigma[u_4]})=4$. The lookahead rule therefore selects
$u_4$, which is not wrong. The remaining six rules are refuted by
explicit stopping games found by exhaustive search over the same family;
each fails on a positive fraction of instances.
\end{proof}

\begin{proposition}[Wrong vertices can be as rare as possible]\label{prop:needle}
For every $k\ge2$ there is a stopping SSG with exactly $k$ non-sink
vertices, all of them Max vertices, and a Max strategy $\sigma$ such that
$S_\sigma=\Vmax$ --- every Max vertex is strictly switchable --- while
exactly one vertex is wrong. Explicitly, let
$\Vmax=\{p,u_1,\dots,u_{k-1}\}$, $\Vmin=\Vavg=\emptyset$, and
\[
  p\to\{t_1,\,t_0\},\qquad u_i\to\{t_1,\,p\}\quad(1\le i\le k-1),
\]
and let $\sigma$ take the second listed successor everywhere, i.e.\
$\sigma(p)=t_0$ and $\sigma(u_i)=p$. Then $|S_\sigma|=k$ and the unique
wrong vertex is $p$. Hence a uniformly random element of $S_\sigma$ is
wrong with probability exactly $1/k$.
\end{proposition}

\begin{proof}
The game is stopping: from $p$ a sink is reached in one step under either
choice, and each $u_i$ moves either to $t_1$ or to $p$, so every play
reaches a sink within two steps.

Its value is $w^\ast\equiv1$: Max plays $p\mapsto t_1$, giving
$w^\ast(p)=1$, and then $w^\ast(u_i)=\max(1,w^\ast(p))=1$. Consequently
$A_p=\{t_1\}$, because $w^\ast(t_1)=1>0=w^\ast(t_0)$, while
$A_{u_i}=\{t_1,p\}$ contains \emph{both} successors, because
$w^\ast(t_1)=1=w^\ast(p)$.

Under $\sigma$ the token goes $u_i\to p\to t_0$, so $\val_\sigma\equiv0$.
Every vertex is therefore strictly switchable: at $p$ we have
$\val_\sigma(t_1)=1>0=\val_\sigma(t_0)$, and at $u_i$ we have
$\val_\sigma(t_1)=1>0=\val_\sigma(p)$. Thus $S_\sigma=\Vmax$ and
$|S_\sigma|=k$. Finally $\sigma(p)=t_0\notin A_p$, so $p$ is wrong, while
$\sigma(u_i)=p\in A_{u_i}$, so no $u_i$ is wrong.
\end{proof}

\begin{remark}
\Cref{prop:needle} is the exact quantitative counterpart of
\Cref{thm:short-path}. That theorem guarantees at least one wrong
switchable vertex; this family shows the guarantee cannot be improved,
since the wrong vertices can be a $1/k$ fraction of $S_\sigma$ with
$k=|\Vmax|$. Choosing a strictly switchable vertex uniformly at random
therefore locates a wrong one with probability only $1/|\Vmax|$ in the
worst case. Note also that every vertex here \emph{is} improving to
switch --- the vertices $u_i$ are not mistakes, they simply fail to reduce
the distance to $\Opt$ --- so no rule can be salvaged by first discarding
non-improving candidates.
\end{remark}

\begin{remark}
\Cref{prop:rules-fail} does not show that these rules require
superpolynomially many iterations; it shows only that the particular
$n$-step argument available from \Cref{thm:short-path} cannot be
completed with them. The distinction matters: the length of improving
paths is provably linear (\Cref{cor:selection}), so any superpolynomial
behaviour must come from a rule repeatedly failing to find the short
path, not from its absence.
\end{remark}

\section{The selection problem is the whole problem}\label{sec:selection}

\Cref{thm:opt-subcube} makes the information needed to solve a stopping
game completely explicit: an optimal strategy is determined by the $n$
bits $\bigl[\,w^\ast(v^{(0)})\ge w^\ast(v^{(1)})\,\bigr]$, one per Max
vertex. We record that these bits are individually exactly as hard as
\Cref{prob:main}, and jointly sufficient.

\begin{definition}
\textsc{SSG-Compare} is the problem: given a \emph{stopping} SSG $G$ and
two vertices $a,b$, decide whether $\val(a)\ge\val(b)$.
\end{definition}

\begin{lemma}[Duality]\label{lem:duality}
Let $G$ be a stopping SSG and let $\bar G$ be obtained from $G$ by
exchanging the two sinks and exchanging the roles of the players, i.e.\
$\Vmax(\bar G)=\Vmin(G)$, $\Vmin(\bar G)=\Vmax(G)$, $\Vavg$ unchanged,
and every edge into $t_1$ redirected to $t_0$ and conversely. Then
$\val_{\bar G}(v)=1-\val_G(v)$ for every \emph{non-sink} $v$.

The restriction to non-sinks is not cosmetic: the two games share the
sinks and their meaning, so $\val_{\bar G}(t_1)=1=\val_G(t_1)$, and the
identity fails there. Every use of the lemma below therefore takes care
that its argument is not a sink.
\end{lemma}

\begin{proof}
Since $G$ is stopping, under every positional pair the token is absorbed
almost surely, so the probabilities of reaching $t_1$ and of reaching
$t_0$ sum to $1$. A pair $(\sigma,\tau)$ of $G$ is the pair
$(\tau,\sigma)$ of $\bar G$ with the roles exchanged, and the $\bar
G$-payoff of that pair at $v$ is $1-v_{\sigma,\tau}(v)$. Hence, using
\Cref{thm:determinacy} twice,
\[
  \val_{\bar G}(v)=\max_\tau\min_\sigma\bigl(1-v_{\sigma,\tau}(v)\bigr)
  =1-\min_\tau\max_\sigma v_{\sigma,\tau}(v)=1-\val_G(v). \qedhere
\]
\end{proof}

\begin{theorem}[Equivalence]\label{thm:compare-equivalence}
\textsc{SSG-Value} and \textsc{SSG-Compare} are polynomial-time
equivalent. Moreover $n$ \emph{non-adaptive} queries to
\textsc{SSG-Compare} suffice to compute an optimal Max strategy, and
hence the exact value vector, of a stopping SSG.
\end{theorem}

\begin{proof}
\emph{\textsc{SSG-Value} $\le$ \textsc{SSG-Compare}.} Let $G$ be an
arbitrary SSG with start vertex $v_0$, and let $G_m$ and
$D=2^{\lceil 3a/2\rceil}$ be as in \Cref{thm:stopping-transform}. One
cannot simply compare $\val_{G_m}(v_0)$ against $\tfrac12$: the
transformation moves the threshold, and comparing at $\tfrac12$ gives the
wrong answer already for the one-vertex game $c\to\{t_1,t_0\}$, where
$\val_G(c)=\tfrac12$ but $\val_{G_m}(c)<\tfrac12$. We must therefore
compare against the shifted threshold, and to do so we manufacture a
vertex whose value is exactly that threshold.

Put
\[
  \theta\;:=\;\tfrac12-\tfrac{3}{8D},
\]
a dyadic rational with denominator $8D=2^{k}$, $k:=\lceil 3a/2\rceil+3$.
Write $\theta=\sum_{i=1}^{k}b_i2^{-i}$ in binary and adjoin to $G_m$ a
chain of $k$ fresh average vertices $c_1,\dots,c_k$ with
\[
  c_i\to\{\,z_i,\;c_{i+1}\,\}\ (i<k),\qquad c_k\to\{z_k,\,t_0\},
  \qquad z_i:=\begin{cases}t_1,&b_i=1\\ t_0,&b_i=0.\end{cases}
\]
Then $\val(c_i)=\tfrac12 [b_i]+\tfrac12\val(c_{i+1})$ with
$\val(c_{k+1}):=0$, so $\val(c_1)=\sum_i b_i2^{-i}=\theta$. The chain
reaches a sink from every $c_i$, no vertex of $G_m$ reaches it, and $k=O(a)$,
so the enlarged game is stopping, has polynomially many vertices, and all
original values are unchanged.

We claim $\val_G(v_0)\ge\tfrac12$ if and only if
$\val_{G_m}(v_0)\ge\val(c_1)$, a single \textsc{SSG-Compare} query. If
$\val_G(v_0)\ge\tfrac12$ then \Cref{thm:stopping-transform} gives
$\val_{G_m}(v_0)>\tfrac12-\tfrac1{4D}=\tfrac12-\tfrac2{8D}>\theta$. If
$\val_G(v_0)<\tfrac12$ then, $\val_G(v_0)$ being a rational of
denominator at most $D$, we get $\val_G(v_0)\le\tfrac12-\tfrac1{2D}$, and
Step~1 of \Cref{thm:stopping-transform} gives
$\val_{G_m}(v_0)\le\val_G(v_0)\le\tfrac12-\tfrac4{8D}<\theta$.

\emph{\textsc{SSG-Compare} $\le$ \textsc{SSG-Value}.} Let $G$ be stopping
with vertices $a,b$. Form the disjoint union of $G$ and its dual
$\bar G$ of \Cref{lem:duality}, sharing the two sinks, and adjoin a fresh
average start vertex $s$ with
\[
  s\;\to\;\{\,a,\ \bar b\,\},
\]
where $\bar b$ is the copy of $b$ in $\bar G$. We may assume $b$ is not a
sink: otherwise replace it by a fresh average vertex $b'$ with both edges
to $b$, which leaves every value unchanged, keeps the game stopping, and
has $\val(b')=\val(b)$. This is exactly the hypothesis
\Cref{lem:duality} requires, and it is not a corner case --- successors
that are sinks are common, and comparing against one directly gives the
wrong answer, since $\val_{\bar G}(t_1)=1$ rather than $0$. The union is
stopping and
has $2N-1$ vertices (the two sinks are shared); the two components do not interact, so their values
are unchanged. By \Cref{lem:duality},
\[
  \val(s)\;=\;\tfrac12\bigl(\val_G(a)+\val_{\bar G}(\bar b)\bigr)
  \;=\;\tfrac12\bigl(1+\val_G(a)-\val_G(b)\bigr),
\]
whence $\val(s)\ge\tfrac12$ if and only if $\val_G(a)\ge\val_G(b)$: a
single \textsc{SSG-Value} query.

\emph{Non-adaptive sufficiency.} Let $G$ be stopping. For each
$v\in\Vmax$ issue the query ``$\val(v^{(0)})\ge\val(v^{(1)})$?''; these
$n$ queries do not depend on one another. Their answers determine, for
each $v$, a successor attaining
$\max\bigl(w^\ast(v^{(0)}),w^\ast(v^{(1)})\bigr)$; let $\sigma$ choose it.
By \Cref{thm:opt-subcube}, $\sigma\in\Opt$. Then $\val_\sigma=w^\ast$ is
computed in polynomial time by the linear program
\begin{equation}\label{eq:lp}
\begin{aligned}
  \text{maximise } &\textstyle\sum_{u\in V} x(u)\quad\text{subject to } x\in[0,1]^V,\ x(t_1)=1,\ x(t_0)=0,\\
  &x(v)=x(\sigma(v)) && (v\in\Vmax),\\
  &x(v)=\tfrac12\bigl(x(v^{(0)})+x(v^{(1)})\bigr) && (v\in\Vavg),\\
  &x(v)\le x(v^{(0)}),\quad x(v)\le x(v^{(1)}) && (v\in\Vmin).
\end{aligned}
\end{equation}
Its feasible set is exactly $\{x\in\mathcal{X}:x\le T_\sigma x\}$, and by
\Cref{cor:comparison} every feasible $x$ satisfies $x\le\val_\sigma$,
while $\val_\sigma$ is itself feasible. Hence $\val_\sigma$ is the unique
\emph{maximiser}, coordinate by coordinate and therefore also for the sum;
the direction matters, since the minimum of the same program is attained
far below $\val_\sigma$.
\end{proof}

The number of bits actually needed can be reduced further: they need not
be indexed by the Max vertices at all, but by the \emph{average} vertices,
and only their relative order matters.

\begin{theorem}[The order of the average values determines the game]\label{thm:order-determines}
Let $G$ be a stopping SSG with $a=|\Vavg|$. The total preorder induced by
$w^\ast$ on $\Vavg\cup\{t_0,t_1\}$ determines $w^\ast$ completely, and
$w^\ast$ is computable from that preorder in polynomial time.
Consequently \textsc{SSG-Value} admits a certificate of $O(a\log a)$ bits,
and $O(a\log a)$ adaptive \textsc{SSG-Compare} queries --- all of them
between average vertices --- suffice to solve the game. The bound involves
neither $|V|$ nor $|\Vmax|$ nor $|\Vmin|$.
\end{theorem}

\begin{proof}
Declare every average vertex a terminal. What remains is a deterministic
max/min game $H$ whose terminal set is $\Vavg\cup\{t_0,t_1\}$, and whose
terminal payoffs, though unknown, are ordered by the given preorder. Let
$K_0\succ K_1\succ\cdots\succ K_\ell$ be the classes, highest first, and
put $Y_i:=K_0\cup\cdots\cup K_i$.

By \Cref{lem:det-game} applied to $H$, the value of a non-terminal $v$ is
the payoff of the best class Max can force, so $v$ belongs to class $K_i$
for the least $i$ with $v\in\mathrm{Attr}(Y_i)$; and if $v$ lies in no
$\mathrm{Attr}(Y_i)$ then Min keeps the play away from every terminal
forever and $v$ has value $0$, the value of $t_0$. Each
$\mathrm{Attr}(Y_i)$ is one backward sweep, so $\ell+1\le a+2$ sweeps
assign every vertex of $G$ a class, using only the preorder.

Every vertex now carries a class, so at each $v\in\Vmax$ we may select a
successor of highest class and at each $v\in\Vmin$ one of lowest class;
let $(\sigma,\tau)$ be any such pair. Because classes order values,
$\sigma$ is greedy for $w^\ast$ and $\tau$ is greedy for Min, whence
$(T_{\sigma,\tau}w^\ast)(v)=(Tw^\ast)(v)=w^\ast(v)$ at every vertex: at
$\Vmax$ and $\Vmin$ by the choice of $\sigma,\tau$, and at $\Vavg$ and the
sinks because the operators agree there. So $w^\ast$ is a fixed point of
$T_{\sigma,\tau}$ and therefore equals $v_{\sigma,\tau}$ by
\Cref{thm:contraction} --- one exact linear solve.

A total preorder on $a+2$ elements is specified by $O(a\log a)$ bits, and
sorting the average vertices by value takes $O(a\log a)$ comparisons, each
an \textsc{SSG-Compare} query between two average vertices.
\end{proof}

\begin{remark}
\Cref{thm:order-determines} strictly sharpens the counting in
\Cref{thm:compare-equivalence} whenever $a<n$: the game is pinned down not
by one bit per Max vertex but by the relative order of the average values,
and the Max and Min vertices contribute nothing beyond what attractor
computations recover for free. Two things it does \emph{not} give. It does
not give a better algorithm --- enumerating preorders costs
$2^{\Theta(a\log a)}$, worse than the $2^{a}\mathrm{poly}(N)$ of
\Cref{thm:few-avg}. And it does not shrink the difficulty: by
\Cref{thm:compare-equivalence} each of the $O(a\log a)$ comparisons is
itself an instance of \Cref{prob:main}. The decoding procedure was
implemented and checked against exhaustive computation on $903$ stopping
games, reproducing the exact value vector every time.
\end{remark}

The atomic step of the selection problem is itself exactly as hard.

\begin{corollary}[Deciding wrongness is the whole problem]\label{cor:wrong-equivalence}
Call $v\in\Vmax$ \emph{wrong} at $\sigma$ if $\sigma(v)\notin A_v$, i.e.\
if $\sigma$ does not choose a $w^\ast$-maximising successor at $v$. The
problem \textsc{Wrong-Vertex} --- given a stopping SSG, a Max strategy
$\sigma$ and a vertex $v\in\Vmax$, decide whether $v$ is wrong at
$\sigma$ --- is polynomial-time equivalent to \textsc{SSG-Value}.
\end{corollary}

\begin{proof}
Write $\{a,b\}$ for the two successors of $v$ with $\sigma(v)=a$. Then
$v$ is wrong at $\sigma$ exactly when $w^\ast(b)>w^\ast(a)$, so a single
\textsc{SSG-Compare} query decides it; by
\Cref{thm:compare-equivalence} this is a polynomial-time reduction to
\textsc{SSG-Value}.

Conversely, given a stopping SSG $H$ and vertices $a,b$, adjoin one fresh
Max vertex $v$ with $v\to\{a,b\}$. Nothing reaches $v$, so all values of
$H$ are unchanged and the enlarged game is still stopping. Let $\sigma$
be any Max strategy of the enlarged game with $\sigma(v)=a$. Then $v$ is
wrong at $\sigma$ if and only if $\val_H(b)>\val_H(a)$, i.e.\ if and only
if the answer to \textsc{SSG-Compare} on $(H,a,b)$ is negative. So
\textsc{SSG-Compare}, and hence \textsc{SSG-Value}, reduces to
\textsc{Wrong-Vertex}.
\end{proof}

\begin{remark}
\Cref{thm:compare-equivalence} sharpens \Cref{cor:selection}. The
difficulty of \textsc{SSG-Value} is neither the length of improving paths
(linear, \Cref{thm:short-path}) nor the amount of information required
(exactly $n$ bits, one per Max vertex, and they need not be obtained
adaptively). It is that each individual bit is itself an instance of the
original problem. Any polynomial algorithm must therefore extract these
$n$ bits by some global mechanism rather than by resolving them one at a
time --- and no such mechanism is supplied by anything proved here.
\end{remark}

\subsection{What a certificate would look like}

It is worth recording exactly what object a polynomial-time algorithm
would have to produce, because the answer is short and shows that the
certificate system is not the difficulty.

\begin{definition}[Bracket]\label{def:bracket}
A \emph{bracket} for a stopping SSG is a pair $(x,y)\in\mathcal{X}^2$ with
$x\le Tx$ and $Ty\le y$. It \emph{decides} an ordered pair of vertices
$(p,q)$ if $x(p)>y(q)$.
\end{definition}

\begin{proposition}[Brackets are sound, complete, and as hard as the problem]\label{prop:bracket}
Let $G$ be a stopping SSG.
\begin{enumerate}[label=(\alph*),leftmargin=2.2em]
\item \emph{(Soundness)} If a bracket decides $(p,q)$ then $w^\ast(p)>w^\ast(q)$.
\item \emph{(Completeness)} $(w^\ast,w^\ast)$ is a bracket, and it decides
  $(p,q)$ whenever $w^\ast(p)>w^\ast(q)$.
\item Deciding whether a given pair $(x,y)$ is a bracket, and whether it
  decides a given $(p,q)$, takes $O(N)$ arithmetic operations.
\item There is a polynomial-time algorithm which, for every stopping SSG
  and every $v\in\Vmax$ whose two successors have distinct values, outputs
  a bracket deciding them in the correct order, \emph{if and only if}
  \textsc{SSG-Value} is decidable in polynomial time.
\end{enumerate}
\end{proposition}

\begin{proof}
(a) By \Cref{cor:comparison}, $x\le Tx$ gives $x\le w^\ast$ and $Ty\le y$
gives $w^\ast\le y$. Hence $w^\ast(p)\ge x(p)>y(q)\ge w^\ast(q)$.
(b) $Tw^\ast=w^\ast$, so $(w^\ast,w^\ast)$ is a bracket, and it decides
exactly the true strict comparisons. (c) Both tests are $N$ evaluations of
$T$ at one coordinate each.

(d) ($\Rightarrow$) Running the algorithm at each $v\in\Vmax$ yields, by
(a), the correct sign of $w^\ast(v^{(0)})-w^\ast(v^{(1)})$ whenever the two
differ, and when it fails to produce a bracket the two values are equal
and either choice is greedy. This is exactly the list of $n$ comparison
bits of \Cref{thm:compare-equivalence}, from which an optimal $\sigma$ and
then $\val$ follow in polynomial time. ($\Leftarrow$) If
\textsc{SSG-Value} is in $\mathrm{P}$ then, by the non-adaptive part of
\Cref{thm:compare-equivalence}, $w^\ast$ is computable in polynomial time,
and $(w^\ast,w^\ast)$ is the required bracket by (b).
\end{proof}

\begin{remark}
So the certificate system is complete and its verification is trivial;
what is missing is a way to \emph{find} a bracket other than by computing
$w^\ast$ itself. \Cref{prop:bracket}(d) makes precise that this is not a
gap that a cleverer certificate format could close --- any sound complete
format has the same property, since by (b) the value vector is always a
witness and by (a) no witness can be wrong. Soundness was also checked
numerically: over $713$ randomly generated brackets on stopping games,
the $3506$ comparisons they decided were correct without exception.
\end{remark}

\section{A subexponential upper bound}\label{sec:randomfacet}

This section contains the strongest positive algorithmic result we can
prove: a Las~Vegas algorithm deciding \Cref{prob:main} in expected time
$e^{2\sqrt n}\,\mathrm{poly}(N)$, where $n=|\Vmax|$. It is subexponential
and not polynomial. Everything is proved; in particular we state and
prove the subgame apparatus that the analysis silently needs.

Throughout this section $G$ is a \emph{stopping} SSG and Max strategies
are identified with $\{0,1\}^{\Vmax}$. Recall $\Phi(x)=\sum_{u\in V}x(u)$
and write $\Phi(\pi):=\Phi(\val_\pi)$.

\subsection{Subcubes, subgames and dead coordinates}

\begin{definition}\label{def:subcube}
Let $F\subseteq\Vmax$ and let $\rho$ assign to each $v\in F$ one of its
successors. The \emph{subcube} $C=C(\rho)$ is the set of Max strategies
extending $\rho$; its \emph{free} coordinates are
$\mathrm{free}(C)=\Vmax\setminus F$ and $n_C:=|\mathrm{free}(C)|$. The
\emph{subgame} $G_C$ is $G$ with \emph{both} out-edges of each $v\in F$
redirected to $\rho(v)$. For a free $w$ and a successor $b$ of $w$,
$C_{w,b}$ denotes the facet of $C$ additionally fixing $w$ to $b$.
\end{definition}

\begin{lemma}[Subgame basics]\label{lem:subgame}
Let $C$ be a subcube of a stopping SSG $G$. Then:
\begin{enumerate}[label=(\alph*),leftmargin=2.2em]
\item $G_C$ is stopping, and the map sending a Max strategy of $G_C$ to its
  restriction-and-completion in $C$ is value-preserving: for every Max
  strategy $\pi'$ of $G_C$ there is $\pi\in C$ inducing the same Markov
  chains, and conversely. (Formally $G_C$ still has $2^{|\Vmax|}$ Max
  strategies, but the choices at the fixed vertices are inert because both
  their out-edges coincide;)
\item for $\pi\in C$ the vector $\val_\pi$ is the same whether computed
  in $G$ or in $G_C$;
\item writing $\opt(C)$ for the value vector of $G_C$, we have
  $\opt(C)\ge\val_\pi$ pointwise for every $\pi\in C$, with equality for
  at least one $\pi\in C$;
\item $C'\subseteq C$ implies $\opt(C')\le\opt(C)$ pointwise;
\item $\pi\in C$ has no strictly switchable coordinate in
  $\mathrm{free}(C)$ if and only if $\val_\pi=\opt(C)$.
\end{enumerate}
\end{lemma}

\begin{proof}
(a) Every positional pair of $G_C$ arises from a positional pair of $G$
(extend $\rho$ arbitrarily) inducing the same Markov chain; since $G$ is
stopping, so is $G_C$. The correspondence of strategies is immediate from
\Cref{def:subcube}. (b) For $\pi\in C$ the chains $M_{\pi,\tau}$ in $G$
and in $G_C$ coincide for every $\tau$, so the pointwise minima over
$\tau$ agree. (c) Apply \Cref{thm:stopping-determinacy} to $G_C$ and use
(b). (d) $G_{C'}$ has fewer Max strategies, all of them strategies of
$G_C$, so the maximum in (c) is over a subset. (e) Apply
\Cref{lem:local-global} to $G_C$, using (b).
\end{proof}

\begin{definition}[Facet potentials and dead coordinates]\label{def:dead}
For a subcube $C$, a free coordinate $w$ and a successor $b$ of $w$ put
$s(C;w,b):=\Phi\bigl(\opt(C_{w,b})\bigr)$. Given also $\sigma\in C$, call
a free coordinate $z$ \emph{dead} for $(C,\sigma)$ if
$s(C;z,\bar\sigma(z))<\Phi(\sigma)$, and let $\mathrm{alive}(C,\sigma)$
be the number of free coordinates that are not dead.
\end{definition}

\begin{lemma}[Facet exclusion]\label{lem:facet-exclusion}
Let $\pi\in C$ and let $w$ be free with successor $b$. If
$\Phi(\pi)>s(C;w,b)$ then $\pi(w)\ne b$.
\end{lemma}

\begin{proof}
If $\pi(w)=b$ then $\pi\in C_{w,b}$, so $\val_\pi\le\opt(C_{w,b})$ by
\Cref{lem:subgame}(c) and hence $\Phi(\pi)\le s(C;w,b)$.
\end{proof}

\subsection{The algorithm and its correctness}

\begin{definition}[Random facet]\label{def:rf}
$\mathrm{RF}(C,\sigma)$, for a subcube $C$ of a stopping SSG and
$\sigma\in C$:
\begin{enumerate}[label=\arabic*.,leftmargin=2.2em]
\item if $n_C=0$, return $(\sigma,\val_\sigma)$;
\item pick $v\in\mathrm{free}(C)$ uniformly at random;
\item $(\sigma^{\!*},\val_{\sigma^{\!*}}):=\mathrm{RF}\bigl(C_{v,\sigma(v)},\sigma\bigr)$;
\item if $v$ is not strictly switchable at $\sigma^{\!*}$ --- the test
  being made in $G$ itself, equivalently in $G_C$, where $v$ is still
  free, and not in $G_{C_{v,\sigma(v)}}$, where its two out-edges have
  been identified --- return
  $(\sigma^{\!*},\val_{\sigma^{\!*}})$;
\item otherwise set $\sigma^{\!**}:=\sigma^{\!*}[v]$ and return
  $\mathrm{RF}\bigl(C_{v,\bar\sigma(v)},\sigma^{\!**}\bigr)$.
\end{enumerate}
A \emph{switch} is an execution of step~5; $S(C,\sigma)$ denotes their
total number.
\end{definition}

\begin{theorem}[Correctness]\label{thm:rf-correct}
$\mathrm{RF}(C,\sigma)$ terminates and returns a strategy $\pi\in C$ with
$\val_\pi=\opt(C)$. Moreover every strategy it holds has value vector
pointwise at least that of its predecessor.
\end{theorem}

\begin{proof}
Induction on $n_C$. If $n_C=0$ then $C=\{\sigma\}$ and
$\val_\sigma=\opt(C)$ by \Cref{lem:subgame}(c). Let $n_C\ge1$ and pick
$v$. By induction $\sigma^{\!*}$ satisfies
$\val_{\sigma^{\!*}}=\opt(C_{v,\sigma(v)})$, and
$\sigma^{\!*}(v)=\sigma(v)$.

If step~4 applies, then $\sigma^{\!*}$ has no strictly switchable
coordinate among $\mathrm{free}(C_{v,\sigma(v)})$, by
\Cref{lem:subgame}(e) applied to $C_{v,\sigma(v)}$, and $v$ is not
strictly switchable either; so $\sigma^{\!*}$ has no strictly switchable
coordinate in $\mathrm{free}(C)$, and \Cref{lem:subgame}(e) applied to
$C$ gives $\val_{\sigma^{\!*}}=\opt(C)$.

Otherwise, by \Cref{lem:switch}, $\val_{\sigma^{\!**}}\ge\val_{\sigma^{\!*}}$
pointwise with strict inequality at $v$, and
$\sigma^{\!**}\in C_{v,\bar\sigma(v)}$. By induction the final call
returns $\pi$ with $\val_\pi=\opt(C_{v,\bar\sigma(v)})\ge\val_{\sigma^{\!**}}$.
Hence $\opt(C_{v,\bar\sigma(v)})\ge\val_{\sigma^{\!*}}=\opt(C_{v,\sigma(v)})$
pointwise and the inequality is strict at $v$. By
\Cref{lem:subgame}(c) some $\pi^{\!*}\in C$ attains $\opt(C)$, and
$\pi^{\!*}$ lies in one of the two facets, so $\opt(C)$ equals
$\opt(C_{v,\sigma(v)})$ or $\opt(C_{v,\bar\sigma(v)})$; since both are
$\le\opt(C)$ by \Cref{lem:subgame}(d) and the second strictly dominates
the first, $\opt(C)=\opt(C_{v,\bar\sigma(v)})=\val_\pi$.

Every switch strictly increases $\Phi$ by \Cref{lem:switch}, and there
are finitely many strategies, so only finitely many switches occur and
the recursion terminates. Monotonicity of the held values is
\Cref{lem:switch} together with the inductive statement.
\end{proof}

\subsection{The expected number of switches}

\begin{definition}\label{def:frec}
Let $f(0):=0$ and, for $m\ge1$,
$f(m):=f(m-1)+1+\frac1m\sum_{j=0}^{m-1}f(j)$.
\end{definition}

Since $f\ge0$, the recursion gives $f(m)-f(m-1)\ge1>0$, so $f$ is
strictly increasing.

\begin{theorem}[Switch bound]\label{thm:rf-bound}
For every stopping SSG, every subcube $C$ and every $\sigma\in C$,
\[
  \mathbb{E}\bigl[S(C,\sigma)\bigr]\;\le\;f\bigl(\mathrm{alive}(C,\sigma)\bigr)\;\le\;f(n_C).
\]
\end{theorem}

\begin{proof}
Strong induction on $n_C$; note that both recursive calls have exactly
$n_C-1$ free coordinates, so the induction is well founded. (The
induction cannot be run on $\mathrm{alive}$, which need not decrease.)
For $n_C=0$ there are no switches and $f(0)=0$.

Write $m:=\mathrm{alive}(C,\sigma)$, $k:=n_C-m$ for the number of dead
coordinates, and
\[
  D\;:=\;\bigl\{w\in\mathrm{free}(C)\;:\;\opt(C_{w,\sigma(w)})\ne\opt(C)\bigr\},
  \qquad d:=|D| .
\]
By \Cref{lem:subgame}(d), $\opt(C_{w,\sigma(w)})\le\opt(C)$ always, so
membership of $D$ means strict pointwise domination.

\smallskip\noindent\emph{Step 1: dead $\Rightarrow$ not in $D$.}
Let $z$ be dead and let $\pi\in C$ attain $\opt(C)$
(\Cref{lem:subgame}(c)). Then $\Phi(\pi)=\Phi(\opt(C))\ge\Phi(\sigma)$.
If $\pi(z)=\bar\sigma(z)$ then \Cref{lem:facet-exclusion} would need
$\Phi(\pi)\le s(C;z,\bar\sigma(z))<\Phi(\sigma)\le\Phi(\pi)$, absurd.
Hence $\pi(z)=\sigma(z)$, so $\pi\in C_{z,\sigma(z)}$ and
$\opt(C_{z,\sigma(z)})\ge\val_\pi=\opt(C)$; with
\Cref{lem:subgame}(d) this forces equality, i.e.\ $z\notin D$. In
particular $D$ consists of alive coordinates and $d\le m$.

\smallskip\noindent\emph{Step 2: a switch occurs if and only if $v\in D$.}
By \Cref{thm:rf-correct}, $\val_{\sigma^{\!*}}=\opt(C_{v,\sigma(v)})$ and
$\sigma^{\!*}(v)=\sigma(v)$. If $v\notin D$ then
$\val_{\sigma^{\!*}}=\opt(C)$, which satisfies the Shapley equation of
$G_C$ at the free Max vertex $v$; hence
$\val_{\sigma^{\!*}}(\bar\sigma(v))\le\val_{\sigma^{\!*}}(v)=\val_{\sigma^{\!*}}(\sigma^{\!*}(v))$
and $v$ is not strictly switchable. If $v\in D$ and $v$ were not strictly
switchable, then as in the proof of \Cref{thm:rf-correct} we would get
$\val_{\sigma^{\!*}}=\opt(C)$, contradicting
$\val_{\sigma^{\!*}}=\opt(C_{v,\sigma(v)})\ne\opt(C)$.

\smallskip\noindent\emph{Step 3: the first recursive call.}
Let $z\ne v$ be dead for $(C,\sigma)$. Since
$C_{v,\sigma(v)}\cap C_{z,\bar\sigma(z)}\subseteq C_{z,\bar\sigma(z)}$,
\Cref{lem:subgame}(d) gives
$s(C_{v,\sigma(v)};z,\bar\sigma(z))\le s(C;z,\bar\sigma(z))<\Phi(\sigma)$,
so $z$ is dead for $(C_{v,\sigma(v)},\sigma)$. Hence the first call has
alive count at most $m-1$ if $v$ is alive, and at most $m$ if $v$ is
dead.

\smallskip\noindent\emph{Step 4: the second recursive call.}
Suppose $v\in D$ is picked and the switch occurs; by Step~1, $v$ is
alive. Since $\sigma\in C_{v,\sigma(v)}$ we have
$\val_\sigma\le\opt(C_{v,\sigma(v)})=\val_{\sigma^{\!*}}$, so
$\Phi(\sigma^{\!*})\ge\Phi(\sigma)$, and by \Cref{lem:switch}
\[
  \Phi(\sigma^{\!**})\;>\;\Phi(\sigma^{\!*})\;=\;s(C;v,\sigma(v))
  \quad\text{and}\quad \Phi(\sigma^{\!**})>\Phi(\sigma).
\]
\emph{(i) The $k$ dead coordinates stay dead.} Let $z\ne v$ be dead for
$(C,\sigma)$. If $\sigma^{\!**}(z)=\bar\sigma(z)$ then
\Cref{lem:facet-exclusion} gives
$\Phi(\sigma^{\!**})\le s(C;z,\bar\sigma(z))<\Phi(\sigma)<\Phi(\sigma^{\!**})$,
absurd; so $\sigma^{\!**}(z)=\sigma(z)$ and therefore
$\overline{\sigma^{\!**}}(z)=\bar\sigma(z)$. Consequently
\[
  s\bigl(C_{v,\bar\sigma(v)};z,\overline{\sigma^{\!**}}(z)\bigr)
  \;=\;s\bigl(C_{v,\bar\sigma(v)};z,\bar\sigma(z)\bigr)
  \;\le\;s(C;z,\bar\sigma(z))\;<\;\Phi(\sigma)\;<\;\Phi(\sigma^{\!**}),
\]
so $z$ is dead for $(C_{v,\bar\sigma(v)},\sigma^{\!**})$.

\emph{(ii) $r(v)-1$ further coordinates become dead.} For $w\in D$ put
$r(w):=\bigl|\{w'\in D: s(C;w',\sigma(w'))\le s(C;w,\sigma(w))\}\bigr|\ge1$.
Let $w\in D$, $w\ne v$, with $s(C;w,\sigma(w))\le s(C;v,\sigma(v))$;
there are $r(v)-1$ such $w$, all alive by Step~1 and all free in
$C_{v,\bar\sigma(v)}$. Then
$\Phi(\sigma^{\!**})>s(C;v,\sigma(v))\ge s(C;w,\sigma(w))$, so
\Cref{lem:facet-exclusion} gives $\sigma^{\!**}(w)\ne\sigma(w)$, i.e.\
$\overline{\sigma^{\!**}}(w)=\sigma(w)$, and hence
\[
  s\bigl(C_{v,\bar\sigma(v)};w,\overline{\sigma^{\!**}}(w)\bigr)
  \;=\;s\bigl(C_{v,\bar\sigma(v)};w,\sigma(w)\bigr)
  \;\le\;s(C;w,\sigma(w))\;<\;\Phi(\sigma^{\!**}),
\]
so $w$ too is dead for $(C_{v,\bar\sigma(v)},\sigma^{\!**})$. The
coordinates in (i) and (ii) are distinct, so the second call has at least
$k+r(v)-1$ dead coordinates among its $n_C-1$ free ones, i.e.\ alive
count at most $(n_C-1)-(k+r(v)-1)=m-r(v)$.

\smallskip\noindent\emph{Step 5: assembling.}
Condition on the uniformly chosen $v$. The bounds of Steps~3 and~4 hold
\emph{deterministically} on the corresponding events, and the randomness
of each recursive call is independent of the choice of $v$, so the
induction hypothesis may be applied inside each branch and the
conditional expectations combined by the tower rule. Using Step~2 for
which branches incur a switch,
\[
  \mathbb{E}[S]\;\le\;\frac1{n_C}\Bigl[\,k\,f(m)\;+\;(m-d)f(m-1)\;+\;\sum_{v\in D}\bigl(f(m-1)+1+f(m-r(v))\bigr)\Bigr].
\]
Ordering $D$ by increasing $s(C;\cdot,\sigma(\cdot))$, its $j$-th element
has $r\ge j$; as $f$ is nondecreasing,
$\sum_{v\in D}f(m-r(v))\le\sum_{j=1}^{d}f(m-j)$. Since $f\ge0$ and
$d\le m$,
\[
  d+\sum_{j=1}^{d}f(m-j)\;\le\;m+\sum_{j=1}^{m}f(m-j)\;=\;m+\sum_{i=0}^{m-1}f(i).
\]
Therefore
\[
  \mathbb{E}[S]\;\le\;\frac1{n_C}\Bigl[k\,f(m)+m\Bigl(f(m-1)+1+\frac1m\sum_{i=0}^{m-1}f(i)\Bigr)\Bigr]
  \;=\;\frac{k\,f(m)+m\,f(m)}{n_C}\;=\;f(m),
\]
by \Cref{def:frec}. If $m=0$ then $D=\emptyset$ by Step~1, no switch ever
occurs, and $\mathbb{E}[S]=0=f(0)$; the term $m\,f(m-1)$ has coefficient
$0$ and the undefined value $f(-1)$ never arises. Finally
$f(m)\le f(n_C)$ because $m\le n_C$ and $f$ is nondecreasing.
\end{proof}

\begin{lemma}\label{lem:fbound}
$f(m)\le e^{2\sqrt m}$ for every $m\ge0$.
\end{lemma}

\begin{proof}
For $m\le8$ this is checked directly: the exact values are
$f(0..8)=0,\,1,\,\tfrac52,\,\tfrac{14}3,\,\tfrac{185}{24},\,\tfrac{713}{60},\,\tfrac{12607}{720},\,\tfrac{62941}{2520},\,\tfrac{1401409}{40320}$,
the largest being $f(8)\approx34.76$ against $e^{2\sqrt8}\approx286.2$.

Let $m\ge9$ and assume the bound for all smaller arguments. Put
$t:=1/\sqrt m\le\tfrac13$. Since $x\mapsto e^{2\sqrt x}$ is increasing,
$\sum_{j=0}^{m-1}e^{2\sqrt j}\le\int_0^{m}e^{2\sqrt x}\,dx=(\sqrt m-\tfrac12)e^{2\sqrt m}+\tfrac12$,
the antiderivative being $(\sqrt x-\tfrac12)e^{2\sqrt x}$. Also
$\sqrt m-\sqrt{m-1}=1/(\sqrt m+\sqrt{m-1})\ge 1/(2\sqrt m)=t/2$, so
$e^{2\sqrt{m-1}}\le e^{2\sqrt m}e^{-t}$. Hence
\[
  f(m)\;\le\;e^{2\sqrt m}e^{-t}+1+\frac1m\Bigl[(\sqrt m-\tfrac12)e^{2\sqrt m}+\tfrac12\Bigr]
  \;=\;e^{2\sqrt m}\bigl[e^{-t}+t-\tfrac{t^{2}}2\bigr]+1+\tfrac{t^{2}}2 .
\]
For $0<t\le1$ the alternating series with decreasing terms gives
$e^{-t}\le1-t+\tfrac{t^2}2-\tfrac{t^3}6+\tfrac{t^4}{24}$, so the bracket
is at most $1-\tfrac{t^3}6+\tfrac{t^4}{24}\le1-\tfrac{t^3}8$. It
therefore suffices that $e^{2\sqrt m}\tfrac{t^3}8\ge1+\tfrac{t^2}2$.
With $u:=\sqrt m\ge3$ this reads $e^{2u}\ge8u^{3}\bigl(1+\tfrac1{2u^{2}}\bigr)$,
which is implied by $e^{2u}\ge12u^{3}$ since
$1+\tfrac1{2u^2}\le\tfrac32$ for $u\ge1$. Finally
$h(u):=2u-\ln12-3\ln u$ satisfies $h(3)=6-\ln12-3\ln3\approx0.219>0$ and
$h'(u)=2-3/u>0$ for $u>\tfrac32$, so $h\ge0$ for all $u\ge3$.
\end{proof}

\begin{theorem}[Subexponential decision procedure]\label{thm:subexp}
There is a Las~Vegas algorithm that decides \Cref{prob:main} --- always
correctly --- in expected time $e^{2\sqrt n}\cdot\mathrm{poly}(N)$, where
$n=|\Vmax|$ and $N=|V|$ of the input game.
\end{theorem}

\begin{proof}
Given $G$ and $v_0$, build the stopping game $G_m$ of
\Cref{thm:stopping-transform}; it has $O(N(a+\log N))$ vertices and the
same set $\Vmax$, hence the same $n$. Run $\mathrm{RF}$ on the full cube
from an arbitrary initial strategy. By \Cref{thm:rf-correct} it returns a
strategy $\pi$ with $\val_\pi=\val_{G_m}$; comparing $\val_\pi(v_0)$
against the threshold $\tfrac12-\tfrac1{4D}$ of
\Cref{thm:stopping-transform} decides $\val_G(v_0)\ge\tfrac12$, so the
output is always correct.

For the running time: by \Cref{thm:rf-bound,lem:fbound} the expected
number of switches is at most $f(n)\le e^{2\sqrt n}$. The recursion tree
has at most $(S+1)(n+1)$ nodes, since between two consecutive switches
the recursion descends a chain of at most $n+1$ calls; by linearity of
expectation the expected number of nodes is $O\bigl(e^{2\sqrt n}n\bigr)$.
Each node performs one evaluation of $\val_\sigma$, which is the unique
optimum of the linear program in the proof of
\Cref{thm:compare-equivalence} --- correct because $G_m$ is stopping, by
\Cref{cor:comparison} --- together with $O(n)$ comparisons of rationals
of polynomial bit-size (\Cref{lem:denominator}). Hence the expected total
time is $e^{2\sqrt n}\cdot\mathrm{poly}(N)$.
\end{proof}

\begin{remark}
The algorithm was implemented and run in exact rational arithmetic on
stopping games; it returned the exact optimum in all $692$ runs, and the
number of switches never exceeded $f(\cdot)$. Independently, the bound of
\Cref{thm:rf-bound} was tested on roughly $13{,}000$ subcube/strategy
pairs drawn from $185$ stopping games with no violation, and each of the
five steps of its proof was checked separately (over $50{,}000$
individual checks). \Cref{lem:fbound} was verified numerically far beyond
the range needed, the ratio $f(m)/e^{2\sqrt m}$ peaking at $0.148$ at
$m=2$ and decreasing thereafter.

We stress what \Cref{thm:subexp} is not. The bound $e^{2\sqrt n}$ is
subexponential but superpolynomial, so this does not touch
\Cref{prob:main}. Nor do we claim the analysis is tight: the recurrence
bound is attained at $\mathrm{alive}=1$, where $\mathbb{E}[S]=f(1)=1$, but
on the instances we could generate the ratio $\mathbb{E}[S]/f(m)$ falls
off steadily with $m$ --- the best observed values were
$0.80,0.71,0.56,0.45$ at $m=2,3,4,5$ --- so the true worst-case behaviour
of the algorithm may well be better than $f(n)$. What we assert is only
that $f(n)$ is a proved upper bound.
\end{remark}

\section{The remaining gap}\label{sec:gap}

Everything above is proved. The following statement is not, and it is
exactly what a polynomial-time algorithm along these lines requires.

\begin{definition}[Switching rule]\label{def:rule}
A \emph{switching rule} is a polynomial-time computable map $R$ that,
given a stopping SSG $G$, the history $\sigma_0,\dots,\sigma_t$ of
strategies produced so far and the value vector $\val_{\sigma_t}$,
returns a nonempty subset $R(\cdot)\subseteq S_{\sigma_t}$ whenever
$S_{\sigma_t}\ne\emptyset$.
\end{definition}

\begin{definition}[The missing statement]\label{def:missing}
\textsc{(Poly-Rule)} \emph{There exist a switching rule $R$ and a
polynomial $p$ such that for every stopping SSG $G$ and every initial
strategy $\sigma_0$, the iteration $\sigma_{t+1}:=\sigma_t[R(\cdot)]$
reaches a strategy with $S_\sigma=\emptyset$ within $p(|V(G)|)$ steps.}
\end{definition}

\begin{theorem}[The gap is exactly the problem]\label{thm:gap-equivalence}
\textsc{(Poly-Rule)} holds if and only if \textsc{SSG-Value} is decidable
in polynomial time.
\end{theorem}

\begin{proof}
($\Rightarrow$) Given $G$ and $v_0$, build $G_m$ by
\Cref{thm:stopping-transform} in polynomial time; $G_m$ is stopping and
of size $O(N(a+\log N))$. Run the iteration of \Cref{def:missing} from an
arbitrary $\sigma_0$. Each step requires evaluating $\val_{\sigma_t}$,
which is the unique optimum of the linear program displayed in the proof
of \Cref{thm:compare-equivalence}; that program's feasible set is
$\{x\in\mathcal{X}:x\le T_{\sigma_t}x\}$ and its optimum is
$\val_{\sigma_t}$ precisely because $G_m$ is stopping, by
\Cref{cor:comparison} --- this is the step that
\Cref{ex:pi-fails} shows cannot be replaced by greedy improvement. Linear
programs of polynomial bit-size are solvable in polynomial time. Then
$S_{\sigma_t}$ is computed by $|\Vmax|$ comparisons. By hypothesis the
iteration halts within $p(|V(G_m)|)$ steps at a strategy with
$S_\sigma=\emptyset$, which by \Cref{lem:local-global} is optimal, so
$\val_{\sigma}=\val_{G_m}$. Comparing $\val_{G_m}(v_0)$ with the explicit
threshold $\tfrac12-\tfrac1{4D}$ of \Cref{thm:stopping-transform} decides
$\val_G(v_0)\ge\tfrac12$. All steps are polynomial.

($\Leftarrow$) Suppose \textsc{SSG-Value} is decidable in polynomial time.
Let $G$ be a stopping SSG. By \Cref{thm:compare-equivalence},
\textsc{SSG-Compare} is then also decidable in polynomial time, so for
each $v\in\Vmax$ we may decide in polynomial time whether
$\val(v^{(0)})\ge\val(v^{(1)})$ and thereby compute, in $n$ queries, a
strategy $\sigma^\ast$ that is greedy with respect to $w^\ast$; by
\Cref{thm:opt-subcube}, $\sigma^\ast\in\Opt$. Define the switching rule
$R$ to be: compute $\sigma^\ast$ as above (this depends only on $G$, not
on the history) and return $\{v\}$ for the least $v\in S_{\sigma_t}$ with
$\sigma_t(v)\ne\sigma^\ast(v)$. The rule is polynomial-time computable,
and such a $v$ exists whenever $S_{\sigma_t}\ne\emptyset$ by
\Cref{thm:short-path}. Each application strictly decreases
$|\{v:\sigma_t(v)\ne\sigma^\ast(v)\}|$, so after at most $n$ iterations
the current strategy either equals $\sigma^\ast$ or has already reached
$S_\sigma=\emptyset$; in both cases the iteration halts within
$n\le|V(G)|$ steps. Thus $R$ witnesses \textsc{(Poly-Rule)} with
$p(x)=x$.
\end{proof}

\begin{remark}
The converse direction is of course not a route to an algorithm --- it
consumes the very oracle one is trying to build. Its content is
negative and precise: \textsc{(Poly-Rule)} is neither weaker nor
stronger than \Cref{prob:main}, so no partial progress on switching rules
can be worth less, or more, than the whole problem. This is why
\Cref{def:missing} is reported as a gap rather than as a plan.
\end{remark}

\begin{remark}[Why the results above do not supply \textsc{(Poly-Rule)}]
\Cref{thm:short-path} guarantees that a path of length at most $n$
exists from every $\sigma$ to $\Opt$, but its proof is
non-constructive in the decisive respect: it selects a disagreement
vertex with a fixed but \emph{unknown} optimal $\sigma^\ast$, and
knowledge of such a $\sigma^\ast$ is precisely knowledge of the answer.
\Cref{thm:all-switches-refuted} rules out the all-switches rule as a
dominating choice, and \Cref{thm:cyclic-uso} rules out potential-based
arguments on the two-player cube. What survives from the analysis is a
guarantee of a different shape: the random-facet rule admits the proved
expected bound $e^{2\sqrt n}$ of \Cref{thm:subexp}, which is
subexponential but not polynomial, and which is tight for that algorithm.
Closing the gap therefore requires a genuinely new mechanism for
\emph{selecting} switches, not a sharper analysis of the rules considered
here.
\end{remark}

\section{Summary}\label{sec:summary}

\paragraph{Foundations, proved from first principles.} The monotone
fixed-point theory of the Shapley operator (\Cref{sec:operators}); for
stopping games, a strategy-independent absorption estimate with exponent
$|\Vavg|$ (\Cref{lem:absorption}), the resulting contraction and
uniqueness of fixed points (\Cref{thm:contraction}), and determinacy
together with the identification of those fixed points with the game
values (\Cref{thm:stopping-determinacy}) --- so that outside
\Cref{thm:determinacy} nothing classical is imported. A denominator bound
$8^{r/2}$ depending only on the number $r$ of average vertices
(\Cref{lem:denominator}), and an explicit reduction to the stopping case
with every constant computed (\Cref{thm:stopping-transform}).

\paragraph{Positive algorithmic results.} Several tractable regimes, each
with a complete proof, and each whose algorithm we implemented and checked
against exhaustive computation:
arbitrary SSGs with $a=O(\log N)$ average vertices, in time
$\mathrm{poly}(N)2^{a}$ and with no stopping hypothesis
(\Cref{thm:few-avg}); games whose cycles are \emph{pure}, that is in which
no cycle passes through an average vertex (\Cref{thm:avg-acyclic}) or,
dually, none passes through a controlled vertex
(\Cref{thm:player-free}), both in polynomial time; one-player games, by
an explicit linear program (\Cref{thm:one-player}); and, in general, a
Las~Vegas algorithm of expected time $e^{2\sqrt n}\,\mathrm{poly}(N)$
obtained from a full analysis of random-facet strategy improvement
(\Cref{thm:rf-correct,thm:rf-bound,lem:fbound,thm:subexp}). The last is
subexponential, not polynomial, and is the strongest general upper bound
here.

\paragraph{Where the difficulty is.} The optimal Max strategies of a
stopping game form exactly the subcube of greedy strategies
(\Cref{thm:opt-subcube}), and from every strategy at most $|\Vmax|$
single strictly improving switches reach it (\Cref{thm:short-path}). So
no lower bound can come from the length of improving paths
(\Cref{cor:selection}); everything is in the \emph{selection} of
switches. That selection problem is exactly the original problem:
\textsc{SSG-Value} is polynomial-time equivalent to comparing two vertex
values, $n$ non-adaptive comparisons determine an optimal strategy
(\Cref{thm:compare-equivalence}), and deciding whether one given vertex
is badly chosen is again equivalent to the whole problem
(\Cref{cor:wrong-equivalence}). The information content is smaller than the
first count suggests --- the order of the average values already determines
everything (\Cref{thm:order-determines}) --- but not the difficulty, since
each comparison is again an instance of the problem.

\paragraph{Refutations and barriers, each with an explicit instance.}
All-switches does not dominate partial switches
(\Cref{thm:all-switches-refuted}) and can take more than $|\Vmax|$
iterations while undoing its own choices
(\Cref{prop:allswitch-overshoot}); the induced orientation of the
combined two-player cube can be cyclic, so no potential lives on it
(\Cref{thm:cyclic-uso}); value iteration needs $2^{\Omega(N)}$ steps
already on games with no players at all (\Cref{thm:vi-lower}); there is
no naive Hamming potential (\Cref{thm:hamming-refuted}) and seven natural
polynomial-time selection rules each fail (\Cref{prop:rules-fail}), while
the wrong vertices can be a $1/|\Vmax|$ fraction of the switchable ones
(\Cref{prop:needle}); the exact gain of a switch is a bounded transfer
impedance divided by an escape probability that can be $2^{-\Theta(N)}$,
so no potential of polynomial range can work (\Cref{thm:impedance}); and
no parameter blind to which player owns which vertex can serve as a
hardness certificate, since it is already large on easy one-player
instances (\Cref{rem:owner-blind}).

\paragraph{The gap.} \Cref{def:missing} states it at full strength and
\Cref{thm:gap-equivalence} proves it equivalent to a polynomial-time
algorithm. What is \emph{not} established anywhere above is a
polynomial-time mechanism for producing even one of the $n$ comparison
bits of \Cref{thm:compare-equivalence} --- equivalently, for selecting a
switchable vertex at which the current strategy is wrong. That is the
whole of the remaining gap, and it is the whole of the problem.

\end{document}
